% Options for packages loaded elsewhere
\PassOptionsToPackage{unicode}{hyperref}
\PassOptionsToPackage{hyphens}{url}
\PassOptionsToPackage{dvipsnames,svgnames,x11names}{xcolor}
\documentclass[
  11pt,
]{article}
\usepackage{xcolor}
\usepackage[a4paper,margin=24mm]{geometry}
\usepackage{amsmath,amssymb}
\setcounter{secnumdepth}{5}
\usepackage{iftex}
\ifPDFTeX
  \usepackage[T1]{fontenc}
  \usepackage[utf8]{inputenc}
  \usepackage{textcomp} % provide euro and other symbols
\else % if luatex or xetex
  \usepackage{unicode-math} % this also loads fontspec
  \defaultfontfeatures{Scale=MatchLowercase}
  \defaultfontfeatures[\rmfamily]{Ligatures=TeX,Scale=1}
\fi
\usepackage{lmodern}
\ifPDFTeX\else
  % xetex/luatex font selection
    \setmainfont[]{TeX Gyre Pagella}
    \setsansfont[]{TeX Gyre Heros}
    \setmonofont[]{DejaVu Sans Mono}
\fi
% Use upquote if available, for straight quotes in verbatim environments
\IfFileExists{upquote.sty}{\usepackage{upquote}}{}
\IfFileExists{microtype.sty}{% use microtype if available
  \usepackage[]{microtype}
  \UseMicrotypeSet[protrusion]{basicmath} % disable protrusion for tt fonts
}{}
\makeatletter
\@ifundefined{KOMAClassName}{% if non-KOMA class
  \IfFileExists{parskip.sty}{%
    \usepackage{parskip}
  }{% else
    \setlength{\parindent}{0pt}
    \setlength{\parskip}{6pt plus 2pt minus 1pt}}
}{% if KOMA class
  \KOMAoptions{parskip=half}}
\makeatother
% definitions for citeproc citations
\NewDocumentCommand\citeproctext{}{}
\NewDocumentCommand\citeproc{mm}{%
  \begingroup\def\citeproctext{#2}\cite{#1}\endgroup}
\makeatletter
 % allow citations to break across lines
 \let\@cite@ofmt\@firstofone
 % avoid brackets around text for \cite:
 \def\@biblabel#1{}
 \def\@cite#1#2{{#1\if@tempswa , #2\fi}}
\makeatother
\newlength{\cslhangindent}
\setlength{\cslhangindent}{1.5em}
\newlength{\csllabelwidth}
\setlength{\csllabelwidth}{3em}
\newenvironment{CSLReferences}[2] % #1 hanging-indent, #2 entry-spacing
 {\begin{list}{}{%
  \setlength{\itemindent}{0pt}
  \setlength{\leftmargin}{0pt}
  \setlength{\parsep}{0pt}
  % turn on hanging indent if param 1 is 1
  \ifodd #1
   \setlength{\leftmargin}{\cslhangindent}
   \setlength{\itemindent}{-1\cslhangindent}
  \fi
  % set entry spacing
  \setlength{\itemsep}{#2\baselineskip}}}
 {\end{list}}
\usepackage{calc}
\newcommand{\CSLBlock}[1]{\hfill\break\parbox[t]{\linewidth}{\strut\ignorespaces#1\strut}}
\newcommand{\CSLLeftMargin}[1]{\parbox[t]{\csllabelwidth}{\strut#1\strut}}
\newcommand{\CSLRightInline}[1]{\parbox[t]{\linewidth - \csllabelwidth}{\strut#1\strut}}
\newcommand{\CSLIndent}[1]{\hspace{\cslhangindent}#1}
\ifLuaTeX
\usepackage[bidi=basic]{babel}
\else
\usepackage[bidi=default]{babel}
\fi
\babelprovide[main,import]{british}
\ifPDFTeX
\else
\babelfont{rm}[]{TeX Gyre Pagella}
\fi
% get rid of language-specific shorthands (see #6817):
\let\LanguageShortHands\languageshorthands
\def\languageshorthands#1{}
\ifLuaTeX
  \usepackage[english]{selnolig} % disable illegal ligatures
\fi
\setlength{\emergencystretch}{3em} % prevent overfull lines
\providecommand{\tightlist}{%
  \setlength{\itemsep}{0pt}\setlength{\parskip}{0pt}}
% Report typography, paired with the markdown-latex-report Lua filter.
% Explicit penalties support the installed LaTeX 2021 kernel.
\usepackage{booktabs,longtable,tabularx,array,colortbl}
\usepackage{xcolor,microtype,fancyhdr,etoolbox,needspace,hyperref,graphicx}
\usepackage{upquote,fvextra}
\usepackage[font=small,labelfont=bf,labelsep=period]{caption}
\definecolor{accent}{HTML}{216B86}
\definecolor{rowalt}{HTML}{F1F6F9}
\definecolor{ruleline}{HTML}{526D7D}
\definecolor{codeborder}{HTML}{CDDCE4}
\arrayrulecolor{ruleline!35}
\setlength{\tabcolsep}{3pt}
\renewcommand{\arraystretch}{1.16}
\setlength{\arrayrulewidth}{0.3pt}
\setlength{\LTpre}{8pt}
\setlength{\LTpost}{8pt}
\clubpenalty=10000
\widowpenalty=10000
\displaywidowpenalty=10000
\brokenpenalty=10000
\setlength{\emergencystretch}{2em}
\hypersetup{colorlinks=true,linkcolor=accent,urlcolor=accent,citecolor=accent}
\makeatletter
\renewcommand\section{\needspace{7\baselineskip}\@startsection{section}{1}{\z@}%
  {-3.2ex plus -1ex minus -.2ex}{1.5ex plus .2ex}%
  {\normalfont\sffamily\Large\bfseries\color{accent}}}
\renewcommand\subsection{\needspace{5\baselineskip}\@startsection{subsection}{2}{\z@}%
  {-2.5ex plus -.5ex minus -.2ex}{1.0ex plus .2ex}%
  {\normalfont\sffamily\large\bfseries\color{accent}}}
\renewcommand\subsubsection{\needspace{4\baselineskip}\@startsection{subsubsection}{3}{\z@}%
  {-2ex plus -.5ex minus -.2ex}{.7ex plus .2ex}%
  {\normalfont\sffamily\normalsize\bfseries}}
\makeatother
\fvset{frame=single,rulecolor=\color{codeborder},framesep=4pt,fontsize=\small,
       breaklines=true,breakanywhere=true,breakindent=1.5em}
\pagestyle{fancy}
\fancyhf{}
\fancyhead[L]{\small\sffamily\color{ruleline}Dough parameter identification}
\fancyhead[R]{\small\sffamily\color{ruleline}\thepage}
\renewcommand{\headrulewidth}{0.3pt}
\setlength{\headheight}{14pt}
\setlength{\headsep}{18pt}
\DeclareGraphicsExtensions{.pdf,.png,.jpg}
\providecommand{\subtwoline}[2]{#1\par\nobreak #2}
\AtBeginEnvironment{CSLReferences}{\small}
\let\reporttableofcontents\tableofcontents
\renewcommand{\tableofcontents}{\reporttableofcontents\clearpage}
\usepackage{bookmark}
\IfFileExists{xurl.sty}{\usepackage{xurl}}{} % add URL line breaks if available
\urlstyle{same}
\hypersetup{
  pdftitle={Dough Parameter Identification},
  pdfauthor={Critical literature survey and independent technical assessment},
  pdflang={en-GB},
  colorlinks=true,
  linkcolor={accent},
  filecolor={Maroon},
  citecolor={accent},
  urlcolor={accent},
  pdfcreator={LaTeX via pandoc}}

\title{Dough Parameter Identification}
\usepackage{etoolbox}
\makeatletter
\providecommand{\subtitle}[1]{% add subtitle to \maketitle
  \apptocmd{\@title}{\par {\large #1 \par}}{}{}
}
\makeatother
\subtitle{Depth recordings, tool trajectories, and the scientific case
for TaichiDough}
\author{Critical literature survey and independent technical assessment}
\date{11 September 2026}

\begin{document}
\maketitle

{
\hypersetup{linkcolor=}
\setcounter{tocdepth}{1}
\tableofcontents
}
\section{Executive assessment}\label{executive-assessment}

\textbf{Your idea is a legitimate scientific research programme:} use
depth recordings and known tool trajectories to calibrate a model of
dough\textquotesingle s elastic, viscous, and plastic response, then ask
whether that model predicts new manipulations. Forward parameter sweeps
are a valid way to perform the inference. A differentiable simulator is
neither a scientific requirement nor a prerequisite for a journal paper.

The broad idea is not new. Visual inverse simulation, robotic material
identification, and inverse dough rheometry each have substantial prior
work. The strongest opportunity is therefore \textbf{not} to claim the
first camera-based dough simulator or the first MPM calibration. It is
to answer a more precise question:

\begin{quote}
Which elastic, rate-dependent, and non-recoverable behaviours can one
calibrated depth camera and known two-tool motion distinguish in real
dough\textemdash{}\allowbreak{}and which inferred parameter combinations
remain predictive outside the calibration motions?
\end{quote}

That question connects robotics, inverse problems, continuum simulation,
and food rheology. A carefully designed study can contribute new
empirical knowledge, a useful calibration method, or a benchmark without
inventing a new MPM solver. The system becomes scientifically persuasive
when its assumptions, limitations, and advantages are measured, not
merely described.

\subsection{The strongest contribution to
pursue}\label{the-strongest-contribution-to-pursue}

The recommended main contribution is an \textbf{observation- and
motion-aware calibration protocol}. It would:

\begin{enumerate}
\def\labelenumi{\arabic{enumi}.}
\tightlist
\item
  distinguish recoverable deformation, rate dependence, and persistent
  deformation through deliberately chosen loading, release, and recovery
  motions;
\item
  estimate only the parameters or parameter combinations that the
  observations constrain;
\item
  account for uncertainty in hidden volume, camera/tool registration,
  contact, and numerical discretisation;
\item
  test continuous predictions on complete, previously unused tool
  trajectories;
\item
  establish what changes when an extra view, an independent mass
  measurement, or a force trace is available.
\end{enumerate}

The novelty would be the demonstrated result of that protocol: for
example, that particular short probe sequences distinguish mechanisms
that arbitrary shaping recordings cannot; that a single view is
sufficient within a quantified operating range; or that apparently
well-fitted material values actually depend on numerical contact
settings. Each is a testable contribution. None is an established result
of TaichiDough yet.

\subsection{What the fresh code assessment
changes}\label{what-the-fresh-code-assessment-changes}

This report uses the working tree at revision \texttt{cb76169} (full
identifier in the evidence appendix), rather than treating the earlier
conversation or report as evidence. It finds a substantive replay,
reconstruction, evaluation, and forward-search implementation. It also
finds three matters that affect interpretation:

\begin{itemize}
\tightlist
\item
  \textbf{The completed material search estimates only an effective
  elastic modulus.} Viscosity is fixed to zero and plastic projection is
  disabled in that experiment. Separate viscosity and contact-retention
  sweeps do not establish joint
  elastic\textendash{}\allowbreak{}viscous\textendash{}\allowbreak{}plastic
  identification.
\item
  \textbf{The previous held-out comparison was wrong.} The
  default-modulus loss of 0.7997 is a \emph{training} score. The
  recorded held-out comparison is 0.8281 for the selected modulus versus
  0.8359 for a frozen-state baseline, approximately 0.93\% lower. There
  is no stored default-modulus held-out score in that result. This small
  improvement has no uncertainty estimate and is not cross-trajectory
  validation.
\item
  \textbf{An affine-transfer scaling inconsistency is present in the
  current solver.} An isolated algebraic test shows that its
  grid-to-particle expression reproduces an affine velocity gradient as
  \(hA\), rather than \(A\), where \(h\) is grid spacing. Because that
  quantity drives deformation, viscous stress, and plastic activation,
  material interpretation requires resolving the inconsistency and
  rerunning calibration. No simulator code was changed for this report.
\end{itemize}

These findings do not make the research idea unscientific. They separate
\emph{verification of the numerical implementation} from
\emph{validation of a model of real dough}. Both are necessary, and they
answer different questions.

\section{Scope, evidence, and how to read this
survey}\label{scope-evidence-and-how-to-read-this-survey}

\subsection{What this review covers}\label{what-this-review-covers}

The review brings together five bodies of work:

\begin{itemize}
\tightlist
\item
  visual inverse simulation and physical parameter identification;
\item
  robotic manipulation and prediction of deformable materials;
\item
  dough rheology and robotic or camera-assisted characterisation;
\item
  MPM discretisation, constitutive modelling, and contact;
\item
  identifiability, model discrepancy, uncertainty, and experimental
  design.
\end{itemize}

A recurring problem is that these literatures use the same words for
different results. A paper may estimate a current shape, identify a
physical parameter, fit a convenient simulator setting, or achieve a
shaping task. All can be valuable, but they cannot be ranked using a
single image error or the presence of automatic differentiation.

\subsection{Review method and limits}\label{review-method-and-limits}

This is a \textbf{critical scoping review and research agenda}, not a
registered systematic review or a proof of priority. Searches used
combinations of dough, viscoelasticity, plasticity, depth/RGB-D, tool
motion, robot manipulation, inverse simulation, system identification,
contact, and parameter calibration. Exact-title follow-up used publisher
records, proceedings, author manuscripts, DOI metadata, arXiv, and open
full-text repositories. Topic evidence files retain search notes, access
limitations, and page/section locations.

Technical claims are strongest when the relevant primary methods and
results were inspected. Where only a publisher abstract or bibliographic
record was available, the report limits itself accordingly. A feature
marked \emph{not established in the inspected material} is not claimed
to be absent from the entire paper or codebase. Reviews are used for
synthesis, not as proof of an uninspected experiment.

The report does not manufacture a screened-paper count, an acceptance
probability, or a claim that no publication has ever combined a
particular set of components. Existing numerical results were inspected,
not independently reproduced by full simulation. The new affine-transfer
check is a small algebraic calculation, explicitly distinguished from a
Taichi rollout. Primary papers downloaded for reading are not bundled
with the web companion.

\subsection{Four different
achievements}\label{four-different-achievements}

\begin{small}
\begin{longtable}{@{}>{\raggedright\sloppy\emergencystretch=3em\arraybackslash}p{0.3003\textwidth} >{\raggedright\sloppy\emergencystretch=3em\arraybackslash}p{0.3015\textwidth} >{\raggedright\sloppy\emergencystretch=3em\arraybackslash}p{0.3232\textwidth}@{}}
\caption{}\\
\toprule
\textbf{Achievement} & \textbf{What it establishes} & \textbf{Evidence it still needs for a stronger claim} \\
\midrule
\endfirsthead
\toprule
\textbf{Achievement} & \textbf{What it establishes} & \textbf{Evidence it still needs for a stronger claim} \\
\midrule
\endhead
Reconstruction or tracking & Agreement with currently observed geometry & Independent future prediction to establish dynamics \\
\hline
Effective simulator calibration & Parameters that improve a specified simulator\textquotesingle s
predictions & Numerical and experimental robustness before physical interpretation \\
\hline
Physical material identification & Mechanically interpretable parameters under stated conditions & Informative loading, uncertainty, and independent mechanical checks \\
\hline
Manipulation or control & Actions that achieve a target or task & Separate identification tests if claiming material recovery \\
\bottomrule
\end{longtable}
\end{small}

A good reconstruction is not a rheological measurement. A broad
parameter uncertainty interval does not necessarily prevent accurate
prediction. A successful controller may compensate for inaccurate
physics through repeated observations. These distinctions are central to
evaluating TaichiDough fairly \citeproc{ref-review2020}{{[}1{]}},
\citeproc{ref-raue2011}{{[}2{]}},
\citeproc{ref-discrepancy2014}{{[}3{]}}.

\section{The inverse problem: what is being
identified?}\label{the-inverse-problem-what-is-being-identified}

\subsection{A measurement-conditioned forward
model}\label{a-measurement-conditioned-forward-model}

Let \(z_t\) denote the simulated material state, \(u_t\) the recorded
tool poses and velocities, \(\theta\) the material parameters, and
\(\psi\) uncertain non-material quantities. A useful formulation is

\[
z_{t+1}=\mathcal S_{h,\Delta t}(z_t,u_t;\theta,\psi),
\qquad
y_t=\mathcal H_{c,o}(z_t)+\delta_t+\epsilon_t.
\]

Here \(h\) and \(\Delta t\) describe numerical resolution;
\(\mathcal H\) projects the state through camera calibration \(c\) and
visibility/occlusion model \(o\); \(\delta\) is model discrepancy; and
\(\epsilon\) is measurement noise. The initial state must also be
reconstructed or measured. The target is not simply a parameter vector
fitted to an arbitrary point cloud: it is a model of the entire
experiment and observation process.

A forward sweep estimates a parameter set by minimising a declared
training objective,

\[
\widehat\theta\in\operatorname*{arg\,min}_{\theta\in\Theta}
\sum_{e\in\mathcal D_{\mathrm{train}}}w_e
\sum_{t\in\mathcal T_e}\ell\!\left(
\mathcal H(\mathcal S^{0:t}(z_{0,e},u_e;\theta,\psi)),y_{t,e}
\right).
\]

The optimiser can be a grid, a coarse-to-fine sweep, random or
space-filling search, evolutionary search, Bayesian optimisation,
finite-difference optimisation, or an adjoint method. These choices
affect computational cost and exploration, not whether the inverse
problem is scientifically meaningful. Finite-difference
Levenberg\textendash{}\allowbreak{}Marquardt is \emph{not autodiff}, but
it is not strictly derivative-free either. This distinction matters when
comparing food-engineering and robotics methods.

\subsection{Separate parameter
families}\label{separate-parameter-families}

\begin{small}
\begin{longtable}{@{}>{\raggedright\sloppy\emergencystretch=3em\arraybackslash}p{0.2321\textwidth} >{\raggedright\sloppy\emergencystretch=3em\arraybackslash}p{0.3403\textwidth} >{\raggedright\sloppy\emergencystretch=3em\arraybackslash}p{0.3525\textwidth}@{}}
\caption{}\\
\toprule
\textbf{Family} & \textbf{Examples} & \textbf{Why the distinction matters} \\
\midrule
\endfirsthead
\toprule
\textbf{Family} & \textbf{Examples} & \textbf{Why the distinction matters} \\
\midrule
\endhead
Bulk material & \(E\), \(\nu\), viscosity \(\eta\), relaxation times, yield/hardening
parameters & These describe a specified constitutive law, not the word ``dough'' in
general \\
\hline
Interface & Tangential slip law, adhesion, floor interaction & Incorrect interfaces can be compensated by apparent bulk stiffness or
dissipation \\
\hline
Geometry and sensing & Hidden volume, mass/density, tool registration, camera pose, timing & These alter inferred strain, force scale, contact onset, and visual
residuals \\
\hline
Numerical & Grid spacing, time step, particle count, transfer, SDF resolution,
contact padding & A fitted value that changes with these settings may describe the
numerical model rather than the material \\
\bottomrule
\end{longtable}
\end{small}

Treating all these quantities as ``dough parameters'' prevents
meaningful physical interpretation. Conversely, holding them fixed does
not remove their uncertainty. It makes every material estimate
conditional on those choices \citeproc{ref-discrepancy2014}{{[}3{]}},
\citeproc{ref-kennedy2001}{{[}4{]}}.

\subsection{A precise meaning of scientific
value}\label{a-precise-meaning-of-scientific-value}

A scientific framework states a hypothesis, makes assumptions explicit,
produces reproducible predictions, and includes experiments that could
contradict its claim. For TaichiDough, a useful claim is not ``the
animation looks realistic.'' It is, for example:

\begin{quote}
Under a specified recipe, sensor configuration, and contact model,
parameters fitted from multi-rate pressing and recovery predict unseen
dual-tool deformation more accurately than parameters fitted from a
single monotonic push, under the same forward-simulation budget.
\end{quote}

Failure to confirm that hypothesis can still provide useful knowledge if
it reveals an observation limit, a constitutive-model failure, or a
dependence on tool contact. An isolated failed software launch, by
contrast, says nothing about rheology.

\section{Dough mechanics and constitutive
interpretation}\label{dough-mechanics-and-constitutive-interpretation}

\subsection{Constitutive models define the
parameters}\label{constitutive-models-define-the-parameters}

Elasticity describes recoverable response. Viscosity describes stress
associated with a deformation rate or flow. Viscoelasticity introduces a
time-dependent combination of elastic and dissipative response, often
with internal state or memory. Plasticity commonly describes
irreversible deformation with a yield condition and loading history.
Food papers sometimes use ``viscoplastic'' more broadly for lasting
viscous flow. A manuscript must give the equation, not rely on the
label.

For a scalar linear Kelvin\textendash{}\allowbreak{}Voigt element,

\[
\sigma=E\varepsilon+\eta\dot\varepsilon.
\]

For a scalar Maxwell element,

\[
\dot\varepsilon=\frac{\dot\sigma}{E}+\frac{\sigma}{\eta}.
\]

They contain similarly named parameters but make different predictions.
At held homogeneous strain,
Kelvin\textendash{}\allowbreak{}Voigt\textquotesingle s rate term
vanishes after the ramp; it does not produce a sustained Maxwell-like
exponential stress relaxation. A standard-linear-solid, generalised
Maxwell model, Burgers model, or fractional/power-law model has a
different response and parameter interpretation. Large-strain dough may
additionally exhibit strain hardening, structural change, rupture, and
interface slip \citeproc{ref-yazar2023}{{[}5{]}},
\citeproc{ref-ng2006}{{[}6{]}}, \citeproc{ref-sun2020}{{[}7{]}}.

The choice need not be the most complicated model. A restricted model
can be scientifically useful if its operating range and failures are
quantified. But adding more sweep values cannot make an incapable
constitutive family reproduce missing physics.

\subsection{Real dough is not
plasticine}\label{real-dough-is-not-plasticine}

Plasticine is a valuable reproducible proxy for robot shaping, but
successful fitting on plasticine is not validation for flour dough.
Flour formulation, hydration, mixing, temperature, resting, and previous
handling affect response. Small-deformation stiffness need not predict
large-deformation behaviour. A sample may become easier or harder to
deform depending on preparation history; experiments themselves change
it \citeproc{ref-yazar2023}{{[}5{]}}.

The food literature already includes nonlinear rheology,
formulation-dependent viscoelastic models, rate-independent hysteretic
descriptions of mixing, and experiments/simulation of bread-dough
rolling \citeproc{ref-ng2006}{{[}6{]}}, \citeproc{ref-sun2020}{{[}7{]}},
\citeproc{ref-anderssen2013}{{[}8{]}},
\citeproc{ref-rolling2009}{{[}9{]}}. These references prevent a claim
that modelling dough with several mechanical parameters is itself new.
Their exact constitutive domains also warn against importing fitted
numbers from extrusion or starch-pasting tests into room-temperature
robot manipulation.

\subsection{Residual deformation does not prove yield
plasticity}\label{residual-deformation-does-not-prove-yield-plasticity}

A remaining deformation after a short wait could result from a yielded
solid, a flowing dashpot, very slow recovery, structural damage,
adhesion to the table, or gravity. A three-second recording cannot
distinguish a permanent shape change from a recovery that takes much
longer without additional evidence.

The fondant study by Cocuzza and Yan makes this issue concrete. It
identifies a three-element spring\textendash{}\allowbreak{}dashpot
model, including a series dashpot, before evaluating gravity and
tool-induced sheet deformation. The series dashpot produces
non-recoverable strain without an explicit yield threshold. The study
adds a different known-force recovery test when tensile data cannot
determine the third parameter reliably
\citeproc{ref-cocuzza2018}{{[}10{]}}. This is a useful precedent for
\emph{changing the experiment to expose a parameter}, rather than
trusting a precise optimiser output.

For TaichiDough, plasticity experiments should include unloading, varied
peak deformation, repeated cycles or matched fresh samples, and
sufficiently long recovery. If the chosen model uses principal-stretch
limits, report those dimensionless limits. Do not rename them a measured
yield stress in pascals.

\subsection{Interfaces are part of the
experiment}\label{interfaces-are-part-of-the-experiment}

Dough can adhere, slip, and stretch during tool separation. Ghorbel and
Launay compare different probe materials and hydration conditions and
report rate-dependent adhesive separation behaviour
\citeproc{ref-adhesion2014}{{[}11{]}}. Their result is not a contact law
for TaichiDough, but it demonstrates why bulk viscosity and tool
adhesion can be confused in a depth sequence.

Pressing, sliding, and withdrawal should be evaluated separately before
a complex two-tool sequence is used to estimate everything
simultaneously. Tool material, moisture, cleanliness, and contact
geometry need documentation. A signed-distance field describes geometry;
it does not select the correct normal, tangential, or adhesive response.

\section{Food-engineering and robotic rheometry
precedents}\label{food-engineering-and-robotic-rheometry-precedents}

\subsection{Inverse semolina-dough
FEM}\label{inverse-semolina-dough-fem}

Fabbri and Cevoli use extrusion measurements and forward FEM to identify
a power-law consistency index \(k\) and flow index \(n\) in
\(\tau=k\dot\gamma^n\) \citeproc{ref-fabbri2015}{{[}12{]}}. These are
flow parameters, not simultaneous elasticity and plasticity estimates.
Known piston/flow motion and die geometry are combined with measured
force-derived pressure.

The inverse method uses Levenberg\textendash{}\allowbreak{}Marquardt
with finite-difference Jacobians. It compares its results against
conventional capillary rheometry with several dies and repeated
measurements. The reported disagreements, 3.6\% for \(k\) and 2.0\% for
\(n\), compare characterisation procedures\textemdash{}\allowbreak{}not
unseen robotic trajectories or camera prediction. The authors do not
claim that agreement identifies which procedure is closer to physical
truth.

\textbf{Relevance:} inverse simulation of dough and numerical parameter
estimation are established. TaichiDough\textquotesingle s possible
distinction is partial visual observation, volumetric tool-conditioned
prediction, and an explicit study of what the measurements identify.
Independent mechanical comparison would substantially strengthen any
claim to physical parameters.

\subsection{MIRANDA and RELAPP}\label{miranda-and-relapp}

Monleón-Getino and colleagues combine computer vision and robotic
material testing, including flour-dough experiments
\citeproc{ref-miranda2025}{{[}13{]}}. MIRANDA tracks visible landmarks
and derives recovery descriptors; RELAPP supplies robotic deformation
and force-related measurements. Regression links descriptors to
laboratory flour-quality quantities such as alveograph strength and RVA
final viscosity.

This is close enough that ``first vision-and-robotics framework for
dough characterisation'' is not a defensible claim. It is also different
from TaichiDough\textquotesingle s proposed result: video recovery
descriptors and regressed industrial indices are not a recovered
three-dimensional continuum material law. RVA final viscosity is not
interchangeable with a room-temperature continuum shear viscosity.

A crucial validation detail is explicit in the full article, §2.5: its
regression data were \textbf{not split into training and test sets}. The
reported fits therefore do not establish withheld-motion or
independent-batch prediction. This creates a useful comparison question,
not a reason to dismiss the work: can a mechanistic,
trajectory-conditioned model make independently tested geometric
predictions from similarly accessible measurements?

\subsection{Rheological robot-shaping and RGB-D
tracking}\label{rheological-robot-shaping-and-rgb-d-tracking}

Cocuzza and Yan identify a rheological model of fondant icing using
tensile and known-force tests, then separately assess sheet deformation
under gravity and a shaping tool \citeproc{ref-cocuzza2018}{{[}10{]}}.
The article predates recent visual MPM systems and shows that robotic
food modelling has an experimental identification history outside MPM.

Petit and colleagues use RGB-D sensing with non-rigid FEM-based tracking
for a pizza-chef application \citeproc{ref-pizza2017}{{[}14{]}}.
Tracking updates the state using current observations. A fitted model
predicting an unseen future trajectory without those updates answers a
different and stronger dynamics question. The proposed TaichiDough paper
should state explicitly whether visual observations initialise,
calibrate, continually correct, or only evaluate the simulation.

Together these studies suggest a productive contribution: connect the
food literature\textquotesingle s controlled preparation and mechanical
tests to robotics\textquotesingle{} partial sensing and
action-conditioned prediction. Hardware uniqueness alone is not the
argument.

\section{Visual and action-conditioned inverse
simulation}\label{visual-and-action-conditioned-inverse-simulation}

\subsection{DPSI: the closest elastoplastic MPM identification
precedent}\label{dpsi-the-closest-elastoplastic-mpm-identification-precedent}

Yang, Ji, and Lai\textquotesingle s DPSI directly addresses
physics-based system identification for robotic manipulation of
elastoplastic materials \citeproc{ref-dpsi2025}{{[}15{]}}. A KUKA robot
executes known tool trajectories on plasticine. A Zivid camera is moved
to six viewpoints before and after each manipulation, and the resulting
point clouds are fused. This is \textbf{one physical camera but
multi-view endpoint observation}, not a continuous fixed-view depth
sequence.

The simulator combines fixed-corotated elasticity, von Mises plasticity,
MLS-MPM/APIC, and rigid-tool SDF contact. Its lower-contact-complexity
experiment fits \(E\), \(\nu\), yield stress, and density; its
higher-contact-complexity experiment also fits table and manipulator
friction. Viscosity is not an explicit parameter in this stated model.
Missing underside geometry is completed under a table-contact
assumption. The tools are coated with flour to reduce adhesion, making
the interface preparation part of the experiment rather than an
incidental detail.

DPSI compares endpoint Chamfer and Earth Mover\textquotesingle s
Distance objectives against observed or volume-filled targets, and uses
Adam. It evaluates a heightmap error separately. The heightmap metric is
a summed residual, not mean depth error in millimetres. Its reported
best parameters are selected using in-distribution validation evaluated
during optimisation; that validation should not be described as an
untouched final test.

Crucially, DPSI already tests six examples of longer, unseen tool
motions, including rolling, repeated poking, and rotation. Those tests
provide genuine action-transfer evidence, although fitted parameters do
not uniformly outperform handpicked ones. The paper also reports flat
loss regions, different low-loss parameter combinations, incomplete
elastic recovery, and contact parameters compensating for material-model
error.

\textbf{Implication for TaichiDough:} known tool motion, incomplete
volumetric reconstruction, elastoplastic MPM fitting, contact
calibration, and unseen-motion evaluation are established. A useful
advance would be to show what continuous depth adds beyond endpoints,
and whether rate-dependent and persistent response can be separated in
\emph{real flour dough}. Merely replacing Adam with a sweep would not
establish that advance.

\subsection{EMPM: already bimanual and already bread
dough}\label{empm-already-bimanual-and-already-bread-dough}

Chen and colleagues\textquotesingle{} EMPM uses three RealSense D455
cameras, action-conditioned MPM, offline human demonstrations, and
online experiments with two Franka arms
\citeproc{ref-empm2026}{{[}16{]}}. Its materials explicitly include
bread dough. It is therefore a direct counterexample to a broad claim
that bimanual, visually calibrated MPM for dough is unexplored.

The formulation includes fixed-corotated elasticity, von Mises
plasticity, and Coulomb table/gripper friction. Its general parameter
notation lists \(E\), \(\nu\), density, and yield stress. However, the
experimental implementation explicitly states that it optimises
homogeneous \textbf{Young\textquotesingle s modulus and
Poisson\textquotesingle s ratio}. A general parameter vector is not
evidence that every entry was jointly identified in the reported
bread-dough experiment. An explicit viscosity parameter is not part of
the listed formulation.

Offline fitting uses point-cloud and tracked-point discrepancies with
AdamW. The paper also explicitly provides \textbf{CMA-ES on the forward
simulator}. This is an especially relevant precedent: derivative-free
embodied MPM identification is already a legitimate published option,
not a scientific deficiency requiring TaichiDough to implement autodiff.

Online fitting omits persistent object tracks and updates near
quasi-static states, comparing short rollouts with observed geometry.
Its bread-dough results measure alignment with and without ongoing
correction. Such adaptation can be useful for robotics, but it is
different from freezing a fitted material law and predicting an
independent deformation rate or new trajectory. The inspected version
does not establish that latter dough result.

\textbf{Implication:} TaichiDough should compare its proposed one-view
measurement requirement, time-resolved loss, explicitly defined
rate/plastic parameters, and frozen-parameter predictions with
EMPM\textquotesingle s actual experiment. Neither ``two arms,'' ``bread
dough,'' nor ``forward optimisation'' alone establishes novelty.

\subsection{DiffCal and DiffCloud: depth is already an identification
input}\label{diffcal-and-diffcloud-depth-is-already-an-identification-input}

Arnavaz and colleagues\textquotesingle{} \emph{Differentiable Depth},
implemented in DiffCal, uses one Intel L515 depth camera with known
silicone-specimen geometry and clamp conditions
\citeproc{ref-diffcal2023}{{[}17{]}}. A stable neo-Hookean FEM with
strain-rate damping is fitted directly against rendered-versus-observed
depth images. Poisson\textquotesingle s ratio is fixed. One dynamic
experiment jointly fits \textbf{Young\textquotesingle s modulus,
damping, and density from twelve depth images at 30 Hz}, following
synchronised support removal.

This is clear prior evidence for single-depth-camera elastic/dissipative
calibration. It is also instructive about physical interpretation: the
paper explicitly treats fitted quantities as parameters of the selected
discrete simulator, and coarse/fine discretisations produce different
stiffness and damping estimates. Its damping coefficient should not be
renamed a dough shear viscosity without matching the equation and units.
Known reference geometry, short dynamic sequences, and largely
fitting-oriented evaluation distinguish it from arbitrary dough with
changing contact and hidden initial volume.

DiffCloud uses two calibrated RealSense cameras, known initial cloth
geometry and grasp location, and recorded robot motion
\citeproc{ref-diffcloud2022}{{[}18{]}}. It fits stiffness and mass
multipliers in a differentiable thin-shell simulator. Its real
experiments commonly minimise a one-way squared point-distance loss at
one selected informative frame rather than supervise every frame.
Folding also uses a visibility-related sampling heuristic.

DiffCloud demonstrates useful real observation alignment. A target
sequence that is subsequently fitted is not an independent
frozen-parameter prediction test, and such a real cross-trajectory test
was not located in the inspected version. Its selected-frame approach is
nevertheless a relevant baseline for TaichiDough: \textbf{does the extra
information in intermediate depth observations actually improve
parameter discrimination and future prediction?} That is an experiment,
not an assumption that more frames must help.

\subsection{Earlier mechanical identification: Wang and
Hahn}\label{earlier-mechanical-identification-wang-and-hahn}

Wang and colleagues combine three Kinect cameras, a high-quality initial
scan, fixed boundaries, and dynamic release under gravity to estimate
soft-object models \citeproc{ref-wang2015}{{[}19{]}}. Corotated FEM uses
\(E\), \(\nu\), and Rayleigh damping parameters \(\alpha\) and
\(\beta\). Their material-fitting subproblem uses \textbf{gradient-free
Nelder\textendash{}\allowbreak{}Mead}, within an alternating procedure
that also estimates tracking and reference geometry. Calling the
material fit derivative-free is accurate; calling every stage
gradient-free would not be.

Separate static loading and dynamic comparisons provide stronger
validation than fitting the original images alone. The paper also
recognises discretisation-dependent apparent stiffness, flat or
multimodal objectives, and the limitations of 30 Hz sampling for
damping. It demonstrates that visual material calibration with
derivative-free optimisation and mechanical validation predates recent
differentiable MPM frameworks.

Hahn and colleagues use ten-camera marker motion capture, known
geometry, and measured clamp motion to identify neo-Hookean elasticity
with power-law viscosity \citeproc{ref-hahn2019}{{[}20{]}}.
Sensitivity/adjoint gradients and L-BFGS fit Lamé parameters, a
viscosity coefficient, and a rate exponent. This is not a markerless
depth method, but it is central to the \emph{mechanics} of the proposed
problem.

Their experiments show why loading diversity matters. Narrow-frequency
motion can poorly constrain a damping law; changing clamp conditions and
combining motions improves estimation; large bending or snapping can
expose constitutive-model limitations. Independent rheometer
measurements use different strain/frequency conditions and are not
automatically exact ground truth for the manipulation experiment.

Together, these studies support a stronger research argument than simply
showing an optimiser converges: demonstrate informative excitation,
disclose numerical dependence, and validate under physically different
conditions.

\subsection{PAC-NeRF and GIC: distinguish synthetic scope from real
evidence}\label{pac-nerf-and-gic-distinguish-synthetic-scope-from-real-evidence}

PAC-NeRF identifies continuum-model parameters from posed multi-view RGB
videos, using differentiable MPM coupled to a radiance-field
representation \citeproc{ref-pacnerf2023}{{[}21{]}}. Its synthetic
studies include elastic, plastic, Newtonian, granular, and non-Newtonian
viscoplastic laws. The last includes shear/bulk response, yield stress,
and plastic viscosity. Therefore a visual inverse problem containing
elastic, viscous, and plastic parameters is already represented in the
literature.

The real demonstration is much narrower: a falling deformable ball
recorded by four cameras, using \textbf{RGB rather than their depth
measurements}. Synthetic recovery under known parameter-generating
models and qualitative real reconstruction should not be described as
experimentally validated joint rheological identification of manipulated
dough.

GIC uses multi-view Gaussian reconstruction to construct geometry and
infer continuum parameters through temporal point-distance and mask
losses \citeproc{ref-gic2024}{{[}22{]}}. Synthetic experiments likewise
cover multiple constitutive families, including viscoplasticity. A
future-frame benchmark tests same-sequence continuation. For sparse-view
real data, the described variant fits masks rather than using the
complete dynamic reconstruction objective. Real grasp demonstrations add
useful application evidence, including transfer of inferred properties
to a different simulation framework, but they are not matched-condition
rheometry.

These systems are strong prior art for visual continuum identification.
They also demonstrate why a survey should separate \textbf{what a
formulation permits}, \textbf{what synthetic experiments recover}, and
\textbf{what real experiments establish}. TaichiDough\textquotesingle s
possible contribution is the latter: an experimentally supported result
about rate, recovery, and residual deformation in real dough under known
tools and limited sensing.

\subsection{Learned constitutive laws and action-conditioned
twins}\label{learned-constitutive-laws-and-action-conditioned-twins}

MASIV learns neural elasticity and plasticity mappings, rather than
recovering a compact vector of \(E\), viscosity, and yield coefficients
\citeproc{ref-masiv2025}{{[}23{]}}. Multi-view reconstructed
trajectories and silhouettes supervise differentiable MPM.
Initialisation and mechanical priors constrain the learned functions;
interior trajectories are inferred rather than directly observed. It is
a relevant alternative when a fixed constitutive family is too
restrictive, but the scientific target differs from interpretable
parameter calibration. Its main-paper evidence includes synthetic
prediction and qualitative transfer, not new real-dough rheometry.

PhysTwin records one- and two-hand interactions with three RGB-D cameras
and fits a spring\textendash{}\allowbreak{}mass model with stiffness,
damping, collision, and control-interaction parameters
\citeproc{ref-phystwin2025}{{[}24{]}}. Zero-order initial fitting is
followed by gradient refinement. A single-image geometry prior does not
make the complete identification method single-camera.

PhysTwin separately reports reconstruction, future-frame continuation,
and unseen interactions. In the latter, the fitted physical model is
reused but registered to the new episode\textquotesingle s initial
observation. This is meaningful transfer evidence with an explicit
initial-state requirement. It shows that bimanual visual fitting and
new-interaction evaluation are established outside MPM as well. The
distinction for TaichiDough must concern the dough law, measurements, or
demonstrated result\textemdash{}\allowbreak{}not just two moving
effectors.

\subsection{Forward and hybrid fitting are established
options}\label{forward-and-hybrid-fitting-are-established-options}

Matl and colleagues identify granular contact coefficients using
depth-derived pile statistics and BayesSim likelihood-free inference
\citeproc{ref-matl2020}{{[}25{]}}. No simulator gradients are required.
Repeated calibration pours at one height are followed by different pour
heights and robotic tasks, with uncertainty and model mismatch examined.
These are granular rather than dough parameters, but the study is a
strong methodological precedent for forward-only calibration aimed at
predictive macroscopic behaviour.

Yoon and Lim combine Bayesian optimisation with later gradient descent
for robot-manipulated cloth \citeproc{ref-yoon2025}{{[}26{]}}. Their
experiments include held-out draping and specimen-size transfer. The
accessible publisher text supports a \textbf{hybrid} method, not a
purely derivative-free one, and comparisons also change simulation
resolution. It should not be cited as isolated evidence that a
particular optimiser is superior.

For TaichiDough, method choice can be practical and transparent: a small
joint grid or space-filling search, selective refinement, and a
plausible-parameter ensemble. Compare its cost and predictive quality
under equal budgets if claiming an inference advantage. A scientific
identification study can instead keep the optimiser conventional and
make the experimental finding its main contribution.

\subsection{Evidence matrix: compare questions, not score
magnitudes}\label{evidence-matrix-compare-questions-not-score-magnitudes}

The following table summarises the primary distinctions. It does not
rank performance across incompatible tasks or declare uninspected
capabilities absent.

\begin{small}
\begin{longtable}{@{}>{\raggedright\sloppy\emergencystretch=3em\arraybackslash}p{0.2547\textwidth} >{\raggedright\sloppy\emergencystretch=3em\arraybackslash}p{0.3221\textwidth} >{\raggedright\sloppy\emergencystretch=3em\arraybackslash}p{0.3481\textwidth}@{}}
\caption{}\\
\toprule
\textbf{Study} & \textbf{What the inspected real experiment identifies or fits} & \textbf{What makes its validation different} \\
\midrule
\endfirsthead
\toprule
\textbf{Study} & \textbf{What the inspected real experiment identifies or fits} & \textbf{What makes its validation different} \\
\midrule
\endhead
DPSI & Plasticine elastic/yield/density and, in a separate contact level,
interface friction; multi-view endpoints & Validation used in parameter selection; additional unseen longer actions \\
\hline
EMPM & Stated experimental \(E,\nu\) fit; three views; bread dough and two arms & Online alignment while adapting is not frozen-parameter dough transfer \\
\hline
DiffCal & Single-view silicone stiffness/damping/density with known geometry & Strong discretisation/fitting study; independent dynamic transfer not
located \\
\hline
DiffCloud & Cloth stiffness/mass multipliers from two views and robot motion & Selected-frame fitting; independent real frozen-parameter transfer not
located \\
\hline
Wang et al. & Elasticity and Rayleigh damping from tracked 3D dynamics & Separate static loading and dynamic checks; derivative-free material fit \\
\hline
Hahn et al. & Elasticity and power-law viscosity from marker motion & Different motions/clamps reveal both information gains and model limits \\
\hline
PAC-NeRF / GIC & Real evidence narrower than synthetic multi-law parameter coverage & Keep rendering, synthetic recovery, temporal continuation, and grasp
tasks separate \\
\hline
PhysTwin & Effective spring/damping/contact model from three-view hand interactions & Future-frame and separate interaction transfer; new initial state
aligned \\
\hline
Fabbri\textendash{}\allowbreak{}Cevoli & Semolina flow consistency and exponent from extrusion pressure & Independent rheometry comparison, not camera-guided action transfer \\
\hline
MIRANDA / RELAPP & Recovery/force descriptors and regressed flour-quality indices & Explicitly no regression train/test split \\
\hline
Proposed TaichiDough study & A defined effective dough law from continuous single-view depth and
recorded tools & Must demonstrate independent prediction and uncertainty; not yet
achieved \\
\bottomrule
\end{longtable}
\end{small}

Three comparisons are particularly valuable for a TaichiDough paper:
\textbf{endpoint versus time-resolved observation}; \textbf{arbitrary
shaping versus designed rate/recovery probes}; and \textbf{fit quality
versus independent frozen-parameter prediction}. These comparisons
provide testable questions that remain meaningful even when the
optimiser and simulator are established methods.

\section{Numerical foundations and what MPM
contributes}\label{numerical-foundations-and-what-mpm-contributes}

\subsection{MLS-MPM, APIC, and a material law are different
choices}\label{mls-mpm-apic-and-a-material-law-are-different-choices}

The original MLS-MPM paper derives a moving-least-squares formulation
and efficient stress-divergence computation; its CPIC extension
addresses displacement discontinuities and two-way rigid coupling
\citeproc{ref-mls2018}{{[}27{]}}. APIC describes the locally affine
velocity representation used in particle\textendash{}\allowbreak{}grid
transfers \citeproc{ref-apic2015}{{[}28{]}}. Neither name specifies the
material\textquotesingle s rheology.

MPM is attractive for large deformation because material state is
carried by particles and computation uses a background grid rather than
a permanently deforming body-fitted mesh. That advantage does not
automatically supply physically correct cutting, fracture, thin-tool
contact, or topology changes. Those require additional formulations and
verification.

The snow MPM model is important because it uses corotated elasticity,
principal-stretch plastic projection, and plastic-volume-dependent
hardening \citeproc{ref-snow2013}{{[}29{]}}. Reusing related mechanisms
for dough can be a useful approximation, but the parameters inherit the
meaning of that particular law. A snow-style dimensionless compression
limit is not a direct rheometer yield-stress measurement.

\subsection{Contact geometry versus contact
mechanics}\label{contact-geometry-versus-contact-mechanics}

An SDF provides signed distance and, where well-defined, a normal. It
can support detecting penetration and projecting velocity or position.
It does not define compliance, restitution, slip, adhesion, or the
relationship between normal and tangential forces.

For example, ILS-MPM explicitly specifies unilateral contact,
normal/tangential penalty treatment, and Coulomb sticking/sliding for
deformable particulate contact \citeproc{ref-ils2020}{{[}30{]}}. Its
evolving level-set formulation is not identical to sampling a static
rigid-tool SDF. It is relevant prior art for contact, but comparison
should name the shared geometric idea and the different mechanical and
numerical treatments.

For physical Coulomb friction, tangential traction is bounded relative
to normal traction. A multiplicative reduction of velocity without
normal-force dependence is a different model. It can be calibrated as an
effective numerical response, but the fitted coefficient must not be
interpreted as a measured Coulomb coefficient.

\subsection{Numerical dissipation can imitate
viscosity}\label{numerical-dissipation-can-imitate-viscosity}

PIC/APIC transfer, grid resolution, time integration, projection, and
repeated contact damping can all affect dissipation. If a per-contact
multiplier \(r\) is applied each time step, its repeated effect scales
like \(r^{T/\Delta t}\) in a simplified continuous-contact example. Its
equivalent decay rate is

\[
\gamma=-\frac{\log r}{\Delta t}.
\]

This is not a constitutive viscosity law. It shows why an unchanged
numerical contact parameter need not represent unchanged dissipation
after a time-step change. In the actual simulator, contact activation
and multiple update stages make the dependence more complicated still.

Likewise, two padding studies answer different questions: holding
padding fixed in metres while refining the grid, or holding
padding/grid-spacing fixed so its physical thickness changes. A paper
should report both quantities and distinguish these experiments.
Material and prediction robustness must be assessed under numerical
changes that preserve mass, initial volume, tool motion, and the
intended physical contact assumptions.

\subsection{Differentiability is an inference
option}\label{differentiability-is-an-inference-option}

ChainQueen establishes differentiable MLS-MPM for soft-robot
applications, while DiffTaichi supplies an efficient differentiable
programming framework \citeproc{ref-chainqueen2019}{{[}31{]}},
\citeproc{ref-difftaichi2020}{{[}32{]}}. These are foundational
references for adjoint-based optimisation, not evidence that every
scientifically useful MPM study must use an adjoint.

A derivative-free search can handle a small number of parameters
transparently and work with discontinuous contact or image operations.
It still suffers from expensive evaluations, parameter coupling, and
ambiguous data. Gradients can accelerate search, but do not create
missing information or guarantee the correct constitutive model. Both
approaches need validation against observations not used for fitting.

\section{Prediction, manipulation, and benchmark
literature}\label{prediction-manipulation-and-benchmark-literature}

\subsection{Learned dynamics is a different scientific
target}\label{learned-dynamics-is-a-different-scientific-target}

RoboCraft reconstructs particle representations from RGB-D, learns graph
dynamics, and uses predictive control for elastoplastic shaping
\citeproc{ref-robocraft2022}{{[}33{]}}. RoboCook adds diverse tools and
long-horizon real dough tasks, with transfer demonstrations to other
materials \citeproc{ref-robocook2023}{{[}34{]}}. \textbf{Real
flour-dough robotic manipulation is therefore not an unexplored
application.}

Their value as comparators depends on the claim. If TaichiDough claims
more accurate forward dynamics, a learned model trained on the same
observations is useful. If it claims interpretable material parameters
or a measurement protocol, an appropriate mechanical/inverse-calibration
comparison may be more important. A complete controller is not mandatory
for a narrowly framed identification paper.

RoboCook\textquotesingle s CEM+MPM versus learned-model/planner
comparisons change more than just calibration. They should not be read
as proof that a carefully calibrated MPM model inherently loses to
learned dynamics. Nor should a TaichiDough static reconstruction IoU be
compared to a final target-shaping IoU as if they measured the same
task.

\subsection{Synthetic and single-view prediction
benchmarks}\label{synthetic-and-single-view-prediction-benchmarks}

PlasticineLab supplies differentiable elastoplastic manipulation tasks
and is useful for controlled synthetic identification experiments
\citeproc{ref-plasticinelab2021}{{[}35{]}}. Recovering parameters from
the exact model that generated the observations is a valuable
implementation test, but it can be an \emph{inverse crime}: the test
omits the model mismatch that makes physical inference difficult. Add
noisy sensing, imperfect registration, altered contact, and ideally a
different observation or forward generator.

DoughNet predicts geometry and topology from a single RGB-D observation
and tool actions \citeproc{ref-doughnet2024}{{[}36{]}}. It does not
establish physical parameter identification in the inspected project
description, but it prevents treating single-view input alone as a new
capability. TaichiDough\textquotesingle s distinction would be the
identifiability, physical interpretation, numerical robustness, or
experimental repeatability it demonstrates.

A useful future benchmark should document more than many depth frames.
It should include independent specimens, preparation metadata,
calibration, tool motion, contact conditions, recovery durations, split
definitions, and reference measurements. This would connect robotics
datasets to food-rheology reproducibility.

\section{What depth and tool motion can
identify}\label{what-depth-and-tool-motion-can-identify}

\subsection{A simple counterexample to absolute-stiffness
identification}\label{a-simple-counterexample-to-absolute-stiffness-identification}

Consider homogeneous small-strain quasistatic elasticity, fixed Poisson
ratio, no body forces, prescribed displacement on part of the boundary,
and zero traction on the remainder:

\[
\nabla\cdot\left[E\,\overline{\mathsf C}(\nu):\varepsilon(u)\right]=0.
\]

For any constant \(a>0\), replacing \(E\) with \(aE\) leaves the
displacement solution unchanged after dividing equilibrium by \(a\).
Reaction forces change, but a depth camera observing displacement does
not measure them. Therefore \textbf{known tool displacement does not, by
itself, guarantee identification of the absolute stiffness scale}.

This is an analytical example with stated assumptions, not a theorem
that all depth-based dough calibration fails. Known gravity, inertia
with independently constrained mass/density, nonzero measured loads, or
force sensing can provide an absolute scale. Whether they provide a
strong enough signal in a particular recording is an experimental
question.

If density is unknown as well, combinations such as \(E/\rho\) can be
easier to infer from dynamics than \(E\) and \(\rho\) separately.
Supplied mass helps, but density computed as mass divided by
reconstructed volume inherits the uncertainty in hidden geometry.

\subsection{A camera may identify timescales better than separate
constants}\label{a-camera-may-identify-timescales-better-than-separate-constants}

For the idealised scalar Kelvin\textendash{}\allowbreak{}Voigt element
during unloaded recovery,

\[
\eta\dot\varepsilon+E\varepsilon=0,
\qquad
\varepsilon(t)=\varepsilon(0)\exp[-tE/\eta].
\]

This recovery identifies the ratio \(\eta/E\) if the assumptions hold.
It does not necessarily identify both absolute parameters. More complex
spatial deformation, inertia, gravity, or additional loads can add
information, but the lesson remains: look for identifiable
\emph{combinations}, not only individual coefficients.

A fixed-displacement hold is especially easy to misinterpret. If the
visible geometry is held still while force relaxes, depth alone cannot
directly measure the stress-relaxation curve. A load cell or independent
force measurement is needed for that direct measurement. Lateral motion
or subsequent release can offer indirect information, but it should not
be described as a measured force-relaxation experiment.

\subsection{Structural identifiability, practical sensitivity, and
prediction}\label{structural-identifiability-practical-sensitivity-and-prediction}

Structural identifiability concerns what ideal observations and the
assumed model can uniquely determine. Practical identifiability concerns
finite, noisy, incomplete measurements. A flat sampled loss curve is
evidence of weak sensitivity under the chosen experiment and objective;
it is not by itself a proof of structural non-identifiability
\citeproc{ref-raue2011}{{[}2{]}}, \citeproc{ref-raue2009}{{[}37{]}}.

There is also \textbf{predictive identifiability}: different parameter
combinations may predict the required output almost identically. A broad
parameter set can therefore support accurate forecasts for one action
class, yet diverge under another. The correct response is to propagate
the plausible set through the intended held-out actions, not
automatically discard the framework because one modulus is uncertain
\citeproc{ref-raue2011}{{[}2{]}}.

Model discrepancy further complicates interpretation. Missing adhesion,
wrong hidden volume, or numerical errors can be absorbed into fitted
mechanical parameters. A narrow numerical optimum is not proof of
physical truth. Explicit discrepancy models can help but also introduce
further confounding; independent information is needed to separate those
effects \citeproc{ref-discrepancy2014}{{[}3{]}},
\citeproc{ref-kennedy2001}{{[}4{]}}.

\subsection{Conditional sweeps are not joint parameter
profiles}\label{conditional-sweeps-are-not-joint-parameter-profiles}

A curve \(L(E,\eta_0,p_0)\) varies elasticity while keeping viscosity
and plasticity fixed. It measures a conditional slice. A profile instead
permits the others to change:

\[
L_{\mathrm{profile}}(E)=\min_{\eta,p,\psi}L(E,\eta,p,\psi).
\]

The latter can reveal compensating combinations hidden by one-at-a-time
sweeps. A joint grid, space-filling search, or repeated conditional
optimisation can approximate a profile without autodiff. A set defined
as ``within 2\% of the best geometric loss'' is a \emph{loss-tolerance
set}, not a 95\% confidence interval. Formal intervals need a defensible
likelihood/noise model or an empirically justified resampling procedure.

With one physical trajectory, resampling the same complete trajectory
cannot estimate between-trajectory uncertainty. Thousands of adjacent
depth frames are not thousands of independent experiments. The sampling
unit must match the claim: trajectory, specimen, batch, or preparation
day.

\subsection{Information-aware motion selection without
gradients}\label{information-aware-motion-selection-without-gradients}

One possible method contribution is to choose a small set of motions
that produces distinguishable outcomes across candidate models. Forward
simulations can estimate feature sensitivities:

\[
J_{jk}\approx\frac{f_j(\theta^{(k,+)})-f_j(\theta^{(k,-)})}
{2\epsilon\,s_j},
\]

where \(\theta^{(k,\pm)}\) equals \(\theta\) except that its positive
\(k\)th parameter is multiplied by \(e^{\pm\epsilon}\). Here \(f_j\) is
a measured trajectory feature and \(s_j\) a noise or scale estimate.
Nearly parallel columns suggest confounded effects; small columns
suggest poor signal. This approximation needs checks across perturbation
sizes and contact regimes. It is a diagnostic, not a guarantee of global
identifiability or a calibrated confidence interval.

Alternatively, no derivative estimate is required: simulate a plausible
parameter ensemble and select motions where their predicted observable
histories disagree most, subject to safe and repeatable contact. Raue
and colleagues provide the general model-based experiment-design
rationale \citeproc{ref-raue2011}{{[}2{]}}. The proposed TaichiDough
contribution would be a concrete real-dough demonstration with matched
acquisition and computation budgets, not a claim to have invented
experiment design.

\section{TaichiDough: current implementation and scientific
interpretation}\label{taichidough-current-implementation-and-scientific-interpretation}

\subsection{What is implemented}\label{what-is-implemented}

The audited code supplies much of the infrastructure needed for the
proposed study:

\begin{itemize}
\tightlist
\item
  metric-camera and floor-plane representations;
\item
  single-view volumetric reconstruction by filling from observed points
  toward the floor;
\item
  particle initialisation that preserves the reconstructed volume and
  supplied mass when changing particle count;
\item
  two-tool replay with timestamp checks, translation interpolation,
  quaternion SLERP, and angular contact-point velocity;
\item
  box, no-tool, and generated-solid SDF collision options;
\item
  visible-depth, mask, coverage, and point-distance evaluation without
  post-hoc per-frame registration or time warping;
\item
  coarse-to-fine modulus search, invalid-candidate handling,
  fingerprints, and declared training/validation windows;
\item
  separate tool-retention and viscosity sweep runners.
\end{itemize}

These are implemented capabilities, not yet a demonstration that all
three material mechanisms are recoverable. Calibration records refer to
metric-v2 transformations, but a schema check is not an independent
metrology error estimate. Generated watertight solids are usable
collision geometry, but their topology tests do not independently verify
physical tool dimensions or marker registration.

\subsection{The actual material
approximation}\label{the-actual-material-approximation}

For deformation state \(F\), polar rotation \(R\), and \(J=\det F\), the
code uses

\[
\tau_e=2\mu(F-R)F^T+\lambda J(J-1)I,
\qquad
\mu=\frac{E}{2(1+\nu)},\quad
\lambda=\frac{E\nu}{(1+\nu)(1-2\nu)}.
\]

This is a corotated stress-like quantity \(PF^T\), not the first Piola
stress itself. When plastic-volume hardening is enabled, both Lamé
coefficients are additionally multiplied by \(\exp[b(1-J_p)]\), where
\(b\) is the hardening setting. The code adds

\[
\tau_v=\eta(C+C^T),
\]

where \(C\) is intended to represent an affine velocity gradient. It is
an instantaneous rate addition, with no generalised Maxwell branch or
relaxation spectrum. At finite deformation, directly adding a rate term
to \(PF^T\) requires a stated stress convention; a spatial Cauchy
viscosity and a Kirchhoff viscosity are not automatically
interchangeable.

Optional plastic projection clips the principal stretches into
\texttt{plastic\_\penalty 250 min}/\texttt{plastic\_\penalty 250 max}.
Optional \texttt{Jp} stores a scalar plastic-volume-related history for
hardening, not a full general plastic state. The limits are
dimensionless. The flag
\texttt{-\penalty 250 -\penalty 250 pure-\penalty 250 viscoelastic}
disables this projection: sweeping plastic bounds while retaining that
flag would provide no evidence about plasticity.

Source locations:
\texttt{scripts/\allowbreak{}taichi\_\penalty 250 viscoelastic\_\penalty 250 mpm\_\penalty 250 scene.py:1268–1303},
\texttt{1397–1407}, and \texttt{2013–2039}. The corresponding reading of
snow-style plasticity should use its explicit stretch/energy equations
rather than assume a yield-stress model
\citeproc{ref-snow2013}{{[}29{]}}.

\subsection{Verified affine-transfer
inconsistency}\label{verified-affine-transfer-inconsistency}

For an interior quadratic B-spline stencil with physical offsets
\(d_i=x_i-x_p\),

\[
\sum_iw_i d_i=0,
\qquad
\sum_iw_i d_i d_i^T=\frac{h^2}{4}I.
\]

An affine grid velocity \(v_i=b+Ad_i\) should be reconstructed by

\[
C=\frac{4}{h^2}\sum_iw_i v_i d_i^T=A.
\]

The audited source instead uses physical
\texttt{dpos=\allowbreak{}(offset-\penalty 250 fx)\allowbreak{}*dx} with
coefficient \texttt{4*inv\_\penalty 250 dx}, which produces
\(C_{\mathrm{code}}=hA\). The coefficient \(4/h\) is appropriate if the
offset is dimensionless; it is inconsistent with the physical offset
used here.

An isolated NumPy reproduction of the exact weights and expression
confirms the issue:

\begin{small}
\begin{longtable}{@{}>{\raggedright\sloppy\emergencystretch=3em\arraybackslash}p{0.2783\textwidth} >{\raggedright\sloppy\emergencystretch=3em\arraybackslash}p{0.2667\textwidth} >{\raggedright\sloppy\emergencystretch=3em\arraybackslash}p{0.3801\textwidth}@{}}
\caption{}\\
\toprule
\textbf{Grid resolution} & \textbf{Input gradient component} & \textbf{Reconstructed component in current expression} \\
\midrule
\endfirsthead
\toprule
\textbf{Grid resolution} & \textbf{Input gradient component} & \textbf{Reconstructed component in current expression} \\
\midrule
\endhead
24 & \(3\ \mathrm{s}^{-1}\) & 0.125 \\
\hline
48 & \(3\ \mathrm{s}^{-1}\) & 0.0625 \\
\hline
96 & \(3\ \mathrm{s}^{-1}\) & 0.03125 \\
\bottomrule
\end{longtable}
\end{small}

The corrected physical-offset formula reproduces \(A\) to floating-point
precision in the same calculation. This is an algebraic verification,
not a measured dough result or a full Taichi regression test. Source:
\texttt{scripts/\allowbreak{}taichi\_\penalty 250 viscoelastic\_\penalty 250 mpm\_\penalty 250 scene.py:1374–1378};
reproducible check:
\texttt{evidence/\allowbreak{}check\_\penalty 250 affine\_\penalty 250 transfer.py}.

Because \(C\) drives the update \((I+\Delta tC)F\), the viscous term,
and plastic activation, the discrepancy affects more than a cosmetic
naming convention. The stored modulus experiment\textquotesingle s
simulator hash matches the audited source. \textbf{Resolve and test this
before interpreting fitted SI parameters or collecting an expensive
calibration dataset.} A correction changes dynamics and requires new
fits; dividing an old fitted modulus by grid resolution is not a
justified substitute. The source code has been left unchanged because
this task is an assessment, not an implementation request.

\subsection{Contact ``friction'' is velocity
retention}\label{contact-friction-is-velocity-retention}

After removing inward normal relative velocity, the tool update
multiplies the remaining relative velocity by

\[
r(1-a)(1-s),
\]

where \(r\) is named
\texttt{tool\_\penalty 250 contact\_\penalty 250 friction}, \(a\) is
absorption, and \(s\) is stickiness. This scales outward normal as well
as tangential relative components and has no normal-traction-based
Coulomb bound. Lower \(r\) means more attenuation; \(r=0\) makes the
corrected particle velocity follow the tool. This direction is not the
interpretation commonly attached to increasing a physical friction
coefficient.

The appropriate current description is \textbf{effective
contact-velocity retention}. Contact geometry, retention, absorption,
stickiness, grid/particle projections, and floor interaction can
compensate for viscosity in a visual loss. Source:
\texttt{scripts/\allowbreak{}taichi\_\penalty 250 viscoelastic\_\penalty 250 mpm\_\penalty 250 scene.py:1335–1343}
and the particle-contact functions. The SDF checks remain useful, but
they validate signed geometry, not a physical friction law.

\subsection{Hidden volume and observation
assumptions}\label{hidden-volume-and-observation-assumptions}

The floor-fill reconstruction assumes that the unobserved volume between
the visible top and floor is occupied. It does not measure hidden
cavities, folded undersides, or detached material. Supplied object mass
divided by this volume determines density. A wrong volume therefore
changes both geometry and mechanical scale. In the inspected modulus
run, the supplied 0.25 kg and reconstructed volume of approximately
\(8.8317\times10^{-4}\) m³ imply about 283 kg/m³. That is a derived
simulation input, not an independent density measurement; the mass and
reconstruction should be checked against the actual specimen.

The current objective averages depth-change, mask-IoU loss,
observed-to-simulated visible-point distance, and real-coverage loss.
These are dimensionless normalised components. Their aggregate is not a
depth error in metres. Dough-only rendering omits tool occlusion, and
the depth-change term uses the intersection of valid initial/current
masks. That can exclude newly exposed or disappearing areas; support and
visibility need separate reporting.

These are reasonable starting assumptions for a restricted task,
provided their effects are measured. Single-view efficiency becomes a
contribution only when the accuracy and uncertainty costs of these
assumptions are quantified.

\subsection{What the stored experiments actually
show}\label{what-the-stored-experiments-actually-show}

The available modulus experiment uses 24,000 particles, grid 48,
\(\Delta t=0.0002\) s, SDF resolution 64, contact padding \(1/384\) m,
supplied mass 0.25 kg, fixed \(\nu=0.3\), zero viscosity, and disabled
plasticity. It selects \(E=130{,}579.320726\) Pa on source frames
1\textendash{}\allowbreak{}60, but reports
\texttt{needs\_\penalty 250 review} and \texttt{accepted: false} because
its sampled near-minimum region is broad.

\begin{small}
\begin{longtable}{@{}>{\raggedright\sloppy\emergencystretch=3em\arraybackslash}p{0.3268\textwidth} >{\raggedright\sloppy\emergencystretch=3em\arraybackslash}p{0.3268\textwidth} >{\raggedright\sloppy\emergencystretch=3em\arraybackslash}p{0.2715\textwidth}@{}}
\caption{}\\
\toprule
\textbf{Setting} & \textbf{Training frames 1\textendash{}\allowbreak{}60} & \textbf{Later frames 61\textendash{}\allowbreak{}97} \\
\midrule
\endfirsthead
\toprule
\textbf{Setting} & \textbf{Training frames 1\textendash{}\allowbreak{}60} & \textbf{Later frames 61\textendash{}\allowbreak{}97} \\
\midrule
\endhead
Selected modulus & 0.361059 & 0.828075 \\
\hline
Default \(E=2000\) Pa & 0.799745 & Not stored in this result \\
\hline
Frozen initial state & 0.343191 & 0.835860 \\
\bottomrule
\end{longtable}
\end{small}

All entries are the same declared aggregate loss; lower is better. The
frozen-state baseline is slightly better on training than the selected
simulation. The selected simulation is about 0.93\% better on the later
window than frozen state, without repetitions or a confidence interval.
These values do not establish a unique modulus or transfer to a new
action. They also \textbf{do not} support the earlier claim that the
fitted model is worse than default \(E\) on held-out data: that claim
incorrectly compared different splits.

Ten sampled modulus values between approximately 91.4 and 186.5 kPa fall
within the configured 2\% loss tolerance. This is weak practical
discrimination for that objective and experiment, not proof of universal
non-identifiability. Its whole-window bootstrap has one training window;
all resamples repeat it, so its zero-width interval provides no
between-trial precision evidence.

The newer \texttt{\_\penalty 250 local} contact-retention sweep
completed, unlike the older failed directory. Its training losses are
0.359191, 0.361067, and 0.361463 for \(r=0.1,0.2,0.3\). The smallest
tested value gives the smallest training score, with only about 0.52\%
improvement relative to 0.2. There is no held-out stage or uncertainty
estimate. This is a completed sensitivity result, not a
physical-friction measurement.

A viscosity runner now tests 0, 1, 2.5, 5, and 10 under the
parameter\textquotesingle s nominal Pa·s interpretation. The available
\texttt{viscosity\_\penalty 250 sweep2} results failed during CUDA
initialisation with
\texttt{CUDA\_\penalty 250 ERROR\_\penalty 250 COMPAT\_\penalty 250 NOT\_\penalty 250 SUPPORTED\_\penalty 250 ON\_\penalty 250 DEVICE}.
They are startup failures, not evidence that the viscosities are
physically wrong or numerically unstable. No successful joint
elasticity\textendash{}\allowbreak{}viscosity\textendash{}\allowbreak{}plasticity
fit was found.

Current metric-v2 static records show a close initial visible-depth
match, but they use the observation from which the initial state was
reconstructed. Older long proxy-tool rollouts are diagnostic evidence
for their own configurations, not a valid measure of current solid-SDF
accuracy. The evidence appendix identifies the files and split
locations.

\section{A falsifiable experimental
programme}\label{a-falsifiable-experimental-programme}

\subsection{Start with the claim, not all available
knobs}\label{start-with-the-claim-not-all-available-knobs}

The first study should target \textbf{predictive calibration of a
restricted material model}, not simultaneous recovery of every
rheological constant. State the preparation, deformation range, loading
speeds, observation duration, tools, and required predictions. Decide
whether the aim is a batch-specific model, transfer between specimens of
one recipe, or transfer between recipes. These are different claims.

A sensible initial model comparison is:

\begin{itemize}
\tightlist
\item
  elastic-only;
\item
  elastic plus a specified rate-dependent term;
\item
  elastic plus a specified plastic mechanism;
\item
  the combined model, if the observations justify its extra parameters.
\end{itemize}

Retain added complexity only when it improves independent prediction or
explains a mechanically verified effect. More parameters and a lower
training loss do not establish better science.

\subsection{Motion primitives and what they
measure}\label{motion-primitives-and-what-they-measure}

\begin{small}
\begin{longtable}{@{}>{\raggedright\sloppy\emergencystretch=3em\arraybackslash}p{0.3027\textwidth} >{\raggedright\sloppy\emergencystretch=3em\arraybackslash}p{0.2950\textwidth} >{\raggedright\sloppy\emergencystretch=3em\arraybackslash}p{0.3272\textwidth}@{}}
\caption{}\\
\toprule
\textbf{Motion or test} & \textbf{Useful information} & \textbf{Key confound or limitation} \\
\midrule
\endfirsthead
\toprule
\textbf{Motion or test} & \textbf{Useful information} & \textbf{Key confound or limitation} \\
\midrule
\endhead
Small press, unload, and release & Reversible response; contact geometry; recovery & Displacement alone may not fix absolute stiffness \\
\hline
Same path at multiple speeds & Rate dependence and candidate timescales & Inertia, contact, timing, and specimen history also change response \\
\hline
Fixed-displacement hold & Force relaxation if force is recorded & Nearly fixed visible geometry may contain little information \\
\hline
Known-force creep and recovery & Compliance, retarded response, lasting flow & A position-controlled robot is not a constant-force apparatus \\
\hline
Load\textendash{}\allowbreak{}unload at several amplitudes & Onset of persistent deformation; model selection & Slow recovery, damage, and adhesion can imitate plasticity \\
\hline
Gravity sag or slump & Known body-force scale and contact-light response & Requires reliable mass/geometry and visible deformation above noise \\
\hline
Controlled slide and withdrawal & Interface slip and adhesion & Contact damping and bulk dissipation remain coupled \\
\hline
New two-tool sequence & Practical action-transfer prediction & Must remain untouched during model and parameter selection \\
\bottomrule
\end{longtable}
\end{small}

Use a pilot to choose speeds, amplitudes, and recovery durations. Values
should resolve the material response above sensor noise while avoiding
unmodelled rupture, uncontrolled adhesion, or tool-tracking errors. A
numerical list chosen merely because its values are convenient is not a
justified physical search range.

One specific warning follows from the current nominal settings:
\(E\approx130.6\) kPa and \(\nu=0.3\) give \(\mu\approx50.2\) kPa.
Ratios \(\eta/\mu\) for \(\eta=1\)\textendash{}\allowbreak{}10 Pa·s are
approximately 20\textendash{}\allowbreak{}200 microseconds, compared
with a 200-microsecond simulation step and roughly 33-millisecond
recorded frame spacing. This is only a dimensional diagnostic; it is not
the exact recovery time of the full model. It does show why the
viscosity range needs an argument based on observable response times,
especially after the affine-transfer inconsistency is resolved.

\subsection{Stage 1: verify the numerical and measurement
chain}\label{stage-1-verify-the-numerical-and-measurement-chain}

Before parameter interpretation:

\begin{enumerate}
\def\labelenumi{\arabic{enumi}.}
\tightlist
\item
  add affine-field reproduction and rigid-motion tests for transfer;
\item
  verify mass/volume invariance across particle counts and grids;
\item
  test simple elastic, rate-dependent, and plastic cases with known
  expected behaviour;
\item
  check camera projection, floor coordinates, tool pose conventions,
  interpolation, and timing independently;
\item
  quantify tool-to-marker and camera-to-scene uncertainty;
\item
  test collision solids and padding without confusing topological
  validity with physical alignment;
\item
  test parameter recovery under the actual depth/visibility operator,
  then deliberately introduce noise and model mismatch.
\end{enumerate}

Report failures separately as launch/environment failures, numerical
failures, missing-observation cases, and physically poor predictions.
Only the last group measures predictive inadequacy. A verification test
does not replace physical validation, but it prevents a search from
compensating for a coding error.

\subsection{Stage 2: fit a small, justified joint parameter
set}\label{stage-2-fit-a-small-justified-joint-parameter-set}

Select bounds from physical response scales, reference tests, and pilot
sensitivity. Positive scale parameters often benefit from logarithmic
sampling; dimensionless stretch bounds require their own physically
meaningful intervals. Represent zero viscosity as an explicit model case
rather than taking its logarithm.

A forward-only protocol can use:

\begin{enumerate}
\def\labelenumi{\arabic{enumi}.}
\tightlist
\item
  a coarse joint grid or space-filling set;
\item
  evaluation of every candidate under identical trajectories, numerical
  settings, and visibility rules;
\item
  refinement around several promising regions, not only the first
  apparent minimum;
\item
  pairwise loss plots and approximate profiles to identify compensation;
\item
  propagation of all plausible parameter sets to validation predictions;
\item
  a final locked evaluation on untouched trials.
\end{enumerate}

A grid with five levels in three parameters has 125 combinations before
repeats, numerical-sensitivity studies, or multiple specimens. This
arithmetic is a compute-planning example, not a proposed universal
budget. Compare sampling strategies under equal forward-evaluation and
physical-data budgets. Record failed evaluations, warm-up/compilation,
simulator time, evaluator time, and total wall time.

One-at-a-time sweeps remain useful diagnostics. They should not be
presented as evidence that elasticity, viscosity, and plasticity were
jointly separated. Similarly, selecting a contact parameter after
examining a former holdout turns that holdout into development data; a
new final test is then needed.

\subsection{Stage 3: separate three kinds of
validation}\label{stage-3-separate-three-kinds-of-validation}

\textbf{Temporal forecast:} fit an early interval and propagate
continuously from the initial state into later frames. Do not reset the
state at the boundary. This tests prediction under the same specimen and
trajectory history.

\textbf{Cross-trajectory prediction:} fit on one set of motions and
predict entirely different recorded motions without refitting. Vary
speed, orientation, indentation, or tool coordination to test the
claimed operating range. This is the central evidence for useful
material calibration.

\textbf{Across-specimen or batch transfer:} test independently prepared
samples or days, with recipe and preparation controlled. A
batch-specific calibration paper does not need to promise universal
recipe transfer, but it must state that limit. If recipe transfer is
claimed, evaluate it explicitly.

An extra camera or a geometric scan can test hidden-volume assumptions.
A force trace or conventional material test can test the mechanical
scale. Those measurements may be reserved for evaluation while the
operational estimator uses one camera. That makes the low-sensor claim
testable rather than assumed.

\subsection{Stage 4: use metrics that reveal different
errors}\label{stage-4-use-metrics-that-reveal-different-errors}

Report camera-space depth error with valid support, silhouette IoU,
bidirectional coverage, and symmetric visible-point distance or a
thresholded F-score. Add geometry features relevant to the task: height,
extent, centroid, spread, or recovered thickness. Force/torque error
belongs in the report when force is available and mechanically relevant.

Always distinguish:

\begin{itemize}
\tightlist
\item
  current-state reconstruction from future prediction;
\item
  visible geometry from complete volume;
\item
  optimised loss from independent evaluation metrics;
\item
  frame-wise curves from episode-level uncertainty;
\item
  a mean error from the coverage on which it was computed.
\end{itemize}

Do not perform per-frame alignment or time warping during final
prediction evaluation. If input registration is uncertain, calibrate it
independently or analyse its uncertainty, rather than aligning away the
prediction error. Include error versus forecast horizon and contact
phase; a final-frame value can conceal wrong intermediate dynamics.

At minimum compare a frozen initial state, unfitted/default mechanics,
and the selected model. For a contact claim, compare calibrated
box/solid-SDF/no-tool cases under otherwise matched settings. For a
multi-parameter claim, compare nested material models. A learned
baseline is important for a learned-versus-physics performance claim,
but not automatically mandatory for a focused identification experiment.

\subsection{Stage 5: uncertainty and numerical
transfer}\label{stage-5-uncertainty-and-numerical-transfer}

Use the physical trajectory, specimen, or batch as the resampling unit.
Choose sample size after a pilot establishes variability and the
smallest practically important paired improvement; this report does not
prescribe an arbitrary number as a journal requirement. Repeats of
interpolation frames do not increase the independent sample count.

Report plausible parameter sets and their predicted outcome spread. Test
camera/tool perturbations, hidden-volume alternatives, contact-retention
assumptions, SDF resolution, particle count, grid spacing, and time
step. Preserve physical mass and geometry when changing resolution.
Separate physical-padding-fixed and grid-relative-padding-fixed tests.

Two outcomes can both be valuable:

\begin{itemize}
\tightlist
\item
  stable parameters and improved independent predictions support
  physical/effective identification within the stated conditions;
\item
  unstable individual parameters but stable useful predictions support a
  narrower predictive calibration claim.
\end{itemize}

If both parameter values and predictions change materially under
plausible numerical or measurement choices, report that limitation. More
printed decimal places do not resolve it.

\section{Where a defensible novel contribution could
lie}\label{where-a-defensible-novel-contribution-could-lie}

\subsection{Novelty is a result, not a list of
components}\label{novelty-is-a-result-not-a-list-of-components}

The existing pieces\textemdash{}\allowbreak{}depth cameras, recorded
tool poses, MPM, SDFs, loss minimisation, and parameter
sweeps\textemdash{}\allowbreak{}are established. Their combination can
still enable a meaningful contribution if it answers a question the
closest studies do not answer in the same experimental conditions. A
narrow hardware conjunction is weak evidence of novelty; a measured
advantage or a carefully established limitation is much stronger.

The following are \textbf{prospective contributions}, ranked here by
their fit to the proposed framework rather than claimed uniqueness in
the entire literature.

\begin{small}
\begin{longtable}{@{}>{\raggedright\sloppy\emergencystretch=3em\arraybackslash}p{0.2776\textwidth} >{\raggedright\sloppy\emergencystretch=3em\arraybackslash}p{0.3398\textwidth} >{\raggedright\sloppy\emergencystretch=3em\arraybackslash}p{0.3076\textwidth}@{}}
\caption{}\\
\toprule
\textbf{Candidate contribution} & \textbf{Testable claim} & \textbf{Evidence that would distinguish it from an implementation report} \\
\midrule
\endfirsthead
\toprule
\textbf{Candidate contribution} & \textbf{Testable claim} & \textbf{Evidence that would distinguish it from an implementation report} \\
\midrule
\endhead
Motion-dependent identification & Designed rate/recovery/amplitude probes distinguish mechanisms better
than arbitrary shaping recordings & Matched-budget comparison, joint parameter sensitivity, and untouched
action-transfer predictions \\
\hline
Single-view measurement sufficiency & One view provides useful predictions within a defined operating range & Independent-view or volume checks; quantify what extra sensing changes \\
\hline
Material/interface separation & Independent contact/timing characterisation reduces false viscosity or
plasticity estimates & Controlled contact perturbations and improved cross-tool or cross-motion
transfer \\
\hline
Numerical robustness of inferred response & Some parameter combinations or predictions remain stable after
refinement & Correct transfer tests and refits across grid/time-step/particle/contact
settings \\
\hline
Real-dough benchmark & Controlled repeated specimens connect rheology to robot prediction & Released preparations, calibrations, depth, trajectories, optional
forces, and locked evaluation splits \\
\bottomrule
\end{longtable}
\end{small}

None requires a new adjoint solver. A more capable constitutive model or
a new search strategy could be an additional contribution, but would
require separate comparison. It is better to demonstrate one principal
contribution clearly than to claim all five on one recording.

\subsection{Recommended primary
hypothesis}\label{recommended-primary-hypothesis}

\textbf{H1:} Multi-rate loading and sufficiently long release/recovery
observations provide more useful constraints on a restricted
elastic\textendash{}\allowbreak{}viscous\textendash{}\allowbreak{}plastic
model than a single monotonic deformation, leading to better prediction
on new two-tool trajectories.

Support would require improvement over a matched-data, matched-compute
baseline, reproducible across independent trials, with the added model
mechanisms actually enabled and numerically verified. H1 is weakened if
the apparent improvement disappears after registration/contact
uncertainty is included, if the inferred response changes with time
step, or if multiple model classes remain observationally equivalent on
new motions.

A complementary hypothesis is that \textbf{not all individual parameters
need be recoverable}: a narrower set of parameter combinations may
suffice for a particular prediction task. Report this as a positive
predictive result with a bounded physical claim, not as complete
rheological recovery.

\subsection{Strong and weak manuscript
claims}\label{strong-and-weak-manuscript-claims}

\textbf{A strong claim, if experimentally supported:}

\begin{quote}
We establish when single-view depth and recorded dual-tool kinematics
are sufficient to calibrate a compact model of real dough response. A
forward-only search over explicitly defined elastic, rate, and
residual-deformation parameters is evaluated on independent tool
trajectories, with contact, hidden-volume, and numerical sensitivity
quantified.
\end{quote}

\textbf{A strong narrower claim, if parameters remain ambiguous:}

\begin{quote}
Several mechanically different parameter sets fit the observed
manipulations. We identify which additional motions or measurements
distinguish them, and which predictions remain reliable despite that
ambiguity.
\end{quote}

\textbf{Claims not supported by the present evidence:} first visual
material identification; first robotic dough rheometry; a novel MLS-MPM
method; physical Coulomb-friction measurement from the current retention
coefficient; joint rheological identification from the completed
elasticity-only run; or universal dough constants.

``An effective parameter for a specified model and preparation'' is not
a concession that makes the work unpublishable. It is an accurate
statement of the inverse problem. It becomes valuable when the resulting
predictions, uncertainty, or limitations are useful to other
researchers.

\section{Journal suitability and the scientific
case}\label{journal-suitability-and-the-scientific-case}

\subsection{Yes, the framework can support a journal
paper}\label{yes-the-framework-can-support-a-journal-paper}

A scientific paper does not have to introduce a new numerical solver.
Official RA-L scope includes innovative ideas, theoretical findings,
application results, and case studies. It does not require gradients or
a new optimiser \citeproc{ref-ral-scope}{{[}38{]}}. RA-P explicitly
accommodates practical investigations, reproducible advances, and case
studies \citeproc{ref-rap-scope}{{[}39{]}}. Research-software
significance can also include substantial engineering rather than
algorithmic novelty, although JOSS has its own maturity, research-use,
licensing, documentation, and testing requirements
\citeproc{ref-joss-scope}{{[}40{]}}.

These scope statements are not acceptance predictions. They show why
``uses known methods'' and ``uses sweeps'' are not sufficient reasons to
reject the idea. Reviewers will still ask what was learned, how it
differs from close prior work, whether experiments test the claims, and
whether another group can reproduce the result.

\subsection{Choose the paper type to match the
evidence}\label{choose-the-paper-type-to-match-the-evidence}

\begin{small}
\begin{longtable}{@{}>{\raggedright\sloppy\emergencystretch=3em\arraybackslash}p{0.2756\textwidth} >{\raggedright\sloppy\emergencystretch=3em\arraybackslash}p{0.3226\textwidth} >{\raggedright\sloppy\emergencystretch=3em\arraybackslash}p{0.3269\textwidth}@{}}
\caption{}\\
\toprule
\textbf{Paper type} & \textbf{Plausible main contribution} & \textbf{What must exist before submission} \\
\midrule
\endfirsthead
\toprule
\textbf{Paper type} & \textbf{Plausible main contribution} & \textbf{What must exist before submission} \\
\midrule
\endhead
Robotics method/experimental paper & Partial-observation calibration with demonstrated action-transfer or
information-aware probing & Verified implementation, closest-method comparison, independent physical
tests, uncertainty \\
\hline
Applied robotics systems paper & Reliable real-dough measurement/replay/calibration procedure with useful
practical findings & Repeated real experiments, specified calibration, failure analysis,
reproducible release \\
\hline
Food-engineering/rheology paper & Mechanically interpretable response under controlled preparation and
loading & Independent force/rheometry reference, correct constitutive
interpretation, batch/repeatability evidence \\
\hline
Simulation/computational-methods paper & New discretisation, contact, constitutive or inference treatment & A genuinely new method plus accuracy/stability/efficiency evaluation \\
\hline
Dataset or scientific-software paper & A reusable calibrated resource or research tool & Data rights/license, documented protocol, working release, reproducible
use beyond one bespoke run \\
\bottomrule
\end{longtable}
\end{small}

For the proposed framework, a focused robotics experimental or
applied-systems article is the most natural direction. A food-rheology
article becomes more plausible with independent mechanical measurements.
A numerical-methods paper requires a numerical contribution; merely
implementing MPM is not one.

A top-level submission should not depend on a long feature
list\textemdash{}\allowbreak{}GUI, network adapters,
reinforcement-learning interfaces\textemdash{}\allowbreak{}unless those
features are used in an experiment. They may belong in documentation.
The paper\textquotesingle s space belongs to the inference problem,
measurements, model, uncertainty, and evidence.

\subsection{What makes this survey itself
useful}\label{what-makes-this-survey-itself-useful}

Broad deformable-object surveys already organise representations,
models, estimation, and control \citeproc{ref-review2020}{{[}1{]}}. Food
reviews already organise dough rheometry and nonlinear behaviour
\citeproc{ref-yazar2023}{{[}5{]}}. A useful survey at their intersection
should organise studies by \textbf{parameter, observation, and
excitation}, rather than by simulator brand or optimiser alone.

This report\textquotesingle s proposed synthesis asks:

\begin{enumerate}
\def\labelenumi{\arabic{enumi}.}
\tightlist
\item
  What physical or effective quantity is estimated?
\item
  What information is measured\textemdash{}\allowbreak{}shape, motion,
  force, or a task outcome?
\item
  What loading and history make that quantity observable?
\item
  Which sensor, contact, geometry, and numerical assumptions remain
  fixed?
\item
  What independent prediction or physical measurement validates the
  result?
\end{enumerate}

That taxonomy makes meaningful comparisons possible between an extrusion
FEM inverse problem, a point-cloud MPM fit, a video-recovery descriptor,
and a learned robot dynamics model. It also makes misleading comparisons
easier to identify. To submit this document itself as a formal review
article would require a more reproducible database/search protocol,
systematic inclusion/exclusion accounting, complete bibliographic
checking, and author-led critical revision. This report is a substantial
starting survey, not a claim that those review-publication requirements
have already been met.

\section{Recommended research
sequence}\label{recommended-research-sequence}

\subsection{Immediate verification and
clarification}\label{immediate-verification-and-clarification}

\begin{enumerate}
\def\labelenumi{\arabic{enumi}.}
\tightlist
\item
  Resolve the affine-transfer inconsistency and add an
  affine-reproduction regression test; rerun affected calibration.
\item
  Define the intended constitutive law and units. Decide whether the
  first paper estimates effective rate and stretch-limit parameters or
  physically interpretable rheological constants.
\item
  Keep elasticity-only, viscosity, plasticity, and combined-model
  experiments distinct. Ensure parameter switches actually activate the
  tested mechanism.
\item
  Rename or clearly document the contact-retention coefficient so it
  cannot be mistaken for Coulomb friction.
\item
  Preserve current logs, exact splits, and source hashes. Do not rewrite
  failed or superseded experiments as successful evidence.
\end{enumerate}

\subsection{First physical study}\label{first-physical-study}

Use controlled pressing/release, multiple speeds, and several amplitudes
on independently prepared specimens. Preserve the depth and pose time
bases, initial mass and volume assumptions, tool calibration, and
preparation history. Use an independent force or geometric reference on
a subset if possible. Compare conditional sweeps with a small joint
search and propagate plausible parameters into unseen motions.

The first result should answer a limited question well: for example,
whether release and recovery observations distinguish elastic-only from
rate-dependent or plastic models better than pressing alone. It need not
recover every parameter simultaneously to be publishable.

\subsection{Full framework study}\label{full-framework-study}

Extend only after the basic identification tests work: cross-trajectory
two-tool prediction, independent specimens/batches, interface
perturbations, numerical refits, and a packaged benchmark. Report both
successful and ambiguous cases. Publish the operational single-camera
pipeline separately from any evaluation-only force/extra-camera
measurements so the claimed sensing requirement remains clear.

\subsection{Suggested reading order}\label{suggested-reading-order}

Read by the question each paper can answer, rather than attempting the
whole bibliography in date order.

\begin{small}
\begin{longtable}{@{}>{\raggedright\sloppy\emergencystretch=3em\arraybackslash}p{0.2073\textwidth} >{\raggedright\sloppy\emergencystretch=3em\arraybackslash}p{0.2555\textwidth} >{\raggedright\sloppy\emergencystretch=3em\arraybackslash}p{0.4621\textwidth}@{}}
\caption{}\\
\toprule
\textbf{Priority} & \textbf{Paper or pair} & \textbf{What to extract for your study} \\
\midrule
\endfirsthead
\toprule
\textbf{Priority} & \textbf{Paper or pair} & \textbf{What to extract for your study} \\
\midrule
\endhead
1 & DPSI \citeproc{ref-dpsi2025}{{[}15{]}} & Closest elastoplastic MPM calibration; contact preparation; endpoint
fitting; unseen actions and failures \\
\hline
2 & EMPM \citeproc{ref-empm2026}{{[}16{]}} & Existing bimanual bread-dough work; actual fitted parameters; adaptive
alignment versus prediction \\
\hline
3 & DiffCal \citeproc{ref-diffcal2023}{{[}17{]}} and Hahn et al.
\citeproc{ref-hahn2019}{{[}20{]}} & What dynamic visual observations constrain; damping definitions; motion
diversity and numerical dependence \\
\hline
4 & Yazar \citeproc{ref-yazar2023}{{[}5{]}} & Dough preparation, nonlinear response, appropriate loading, recovery and
interface conditions \\
\hline
5 & Fabbri\textendash{}\allowbreak{}Cevoli
\citeproc{ref-fabbri2015}{{[}12{]}} and
Cocuzza\textendash{}\allowbreak{}Yan
\citeproc{ref-cocuzza2018}{{[}10{]}} & Independent mechanical checks; change the experiment when a parameter
remains unconstrained \\
\hline
6 & Raue et al. \citeproc{ref-raue2011}{{[}2{]}} & Conditional slices versus profiles; uncertain parameters versus reliable
predictions \\
\hline
7 & MLS-MPM and snow MPM \citeproc{ref-mls2018}{{[}27{]}},
\citeproc{ref-snow2013}{{[}29{]}} & Transfer/stress equations and the precise meaning of stretch-based
plasticity \\
\hline
8 & PAC-NeRF and GIC \citeproc{ref-pacnerf2023}{{[}21{]}},
\citeproc{ref-gic2024}{{[}22{]}} & Joint viscoplastic parameterisation; synthetic versus real evidence \\
\hline
9 & Matl et al. \citeproc{ref-matl2020}{{[}25{]}} and Wang et al.
\citeproc{ref-wang2015}{{[}19{]}} & Forward-only material fitting, uncertainty and independent validation \\
\hline
10 & MIRANDA/RELAPP, RoboCook and PhysTwin
\citeproc{ref-miranda2025}{{[}13{]}},
\citeproc{ref-phystwin2025}{{[}24{]}},
\citeproc{ref-robocook2023}{{[}34{]}} & Adjacent food measurement, practical shaping, and new-interaction
prediction \\
\bottomrule
\end{longtable}
\end{small}

The reading catalogue records a primary link, the source version
inspected, and whether access was full, partial, abstract-level, or
metadata-only. That distinction is especially important before copying
equations, reported parameter values, or evaluation claims into a
manuscript.

\section{Evidence appendix}\label{evidence-appendix}

\subsection{Current repository
records}\label{current-repository-records}

The audit is tied to revision
\texttt{cb76169436d334c127f165ce5c3f779100bb82d4}. Relative paths below
avoid dependence on a particular machine\textquotesingle s home
directory.

\begin{small}
\begin{longtable}{@{}>{\raggedright\sloppy\emergencystretch=3em\arraybackslash}p{0.4931\textwidth} >{\raggedright\sloppy\emergencystretch=3em\arraybackslash}p{0.4444\textwidth}@{}}
\caption{}\\
\toprule
\textbf{Record} & \textbf{What was checked} \\
\midrule
\endfirsthead
\toprule
\textbf{Record} & \textbf{What was checked} \\
\midrule
\endhead
\texttt{scripts/\allowbreak{}taichi\_\penalty 250 viscoelastic\_\penalty 250 mpm\_\penalty 250 scene.py:1268–1407} & Material terms, plastic switches, transfers, and contact/velocity
updates \\
\hline
\texttt{scripts/\allowbreak{}deformpath\_\penalty 250 dynamics.py:390–516} & Timestamped two-tool interpolation and velocities \\
\hline
\texttt{scripts/\allowbreak{}reconstruct\_\penalty 250 voxel\_\penalty 250 dough\_\penalty 250 from\_\penalty 250 deformpath.py:223–372} & Floor-fill and hidden-volume assumptions \\
\hline
\texttt{scripts/\allowbreak{}material\_\penalty 250 calibration.py:164–332} & Normalised loss, support, and validity rules \\
\hline
\texttt{scripts/\allowbreak{}calibrate\_\penalty 250 youngs\_\penalty 250 modulus.py:1184–1379} & Search, selection, split handling, and reporting \\
\hline
\texttt{data/\allowbreak{}single\_\penalty 250 episode\_\penalty 250 calibration/\allowbreak{}episode18\_\penalty 250 kugla\_\penalty 250 fixed\_\penalty 250 collision\_\penalty 250 11\_\penalty 250 9/\allowbreak{}material\_\penalty 250 calibration\_\penalty 250 result.json} & Correct training/held-out/default/frozen-state comparisons \\
\hline
\texttt{output/\allowbreak{}episode18\_\penalty 250 e130579\_\penalty 250 dx8\_\penalty 250 tool\_\penalty 250 friction\_\penalty 250 sweep\_\penalty 250 local/\allowbreak{}friction\_\penalty 250 sweep.json} & Completed training-only retention sweep \\
\hline
\texttt{output/\allowbreak{}episode18\_\penalty 250 e130579\_\penalty 250 dx8\_\penalty 250 viscosity\_\penalty 250 sweep2/\allowbreak{}viscosity\_\penalty 250 sweep.json} & Failed startup cases; stderr identifies CUDA initialisation failure \\
\bottomrule
\end{longtable}
\end{small}

In the large material-result JSON, selection/flat-minimum/one-window
bootstrap are at lines 227\textendash{}\allowbreak{}289; the held-out
aggregate is at 86830\textendash{}\allowbreak{}86837; the default
comparison explicitly carries the training split at
93862\textendash{}\allowbreak{}93899. The full local evidence note
contains additional locations and interpretation. Numeric reports under
\texttt{data/\allowbreak{}} and external recorded inputs require an
archive or download procedure; source code alone is not a complete
reproduction package.

\subsection{What was and was not
executed}\label{what-was-and-was-not-executed}

Executed for this assessment: read-only source/result inspection,
primary-literature retrieval, bibliography preparation, the isolated
affine-weight calculation, and document builds. \textbf{Not executed:} a
new material calibration, a new robot experiment, a full Taichi
regression suite, or an independent reproduction of the stored real-data
results. No simulator or sweep source files were modified.

The accompanying evidence files distinguish full-text inspection,
abstract-level evidence, and metadata-only reading leads. The
bibliography identifies the particular editions cited; related
conference/journal titles and online/issue years may differ. The older
report remains a separate file and should not be used for its erroneous
mixed-split comparison.

\subsection{Literature access and version
qualifications}\label{literature-access-and-version-qualifications}

\begin{small}
\begin{longtable}{@{}>{\raggedright\sloppy\emergencystretch=3em\arraybackslash}p{0.4459\textwidth} >{\raggedright\sloppy\emergencystretch=3em\arraybackslash}p{0.4916\textwidth}@{}}
\caption{}\\
\toprule
\textbf{Source group} & \textbf{Scientific text inspected and limit} \\
\midrule
\endfirsthead
\toprule
\textbf{Source group} & \textbf{Scientific text inspected and limit} \\
\midrule
\endhead
DPSI; EMPM & Full arXiv v3 (February \penalty 5000 2025) and v1 (January \penalty 5000 2026), respectively; final
journal publication metadata checked separately \\
\hline
PAC-NeRF; GIC; DiffCloud & Full arXiv versions v1 (2023), v3 (2024), and v2 (2025);
DiffCloud\textquotesingle s conference publication is 2022 \\
\hline
PhysTwin; MASIV & CVF accepted main papers; separate supplements not read \\
\hline
DiffCal; Cocuzza\textendash{}\allowbreak{}Yan & Full early-online/accepted author manuscripts; final issue metadata
checked separately \\
\hline
Wang; Hahn; Matl; Fabbri; MIRANDA; Yazar; RoboCook & Full papers or relevant primary full-text methods/results inspected \\
\hline
MLS-MPM; snow MPM; ILS-MPM; Raue 2011; Arriola-Rios review & Relevant primary full-text numerical/methodological sections inspected \\
\hline
APIC; RoboCraft; PlasticineLab; DoughNet; ChainQueen; DiffTaichi & Primary abstracts/project descriptions and publication records; detailed
results not inferred from those alone \\
\hline
Raue 2009; Kennedy\textendash{}\allowbreak{}O\textquotesingle Hagan;
Brynjarsdóttir\textendash{}\allowbreak{}O\textquotesingle Hagan;
Anderssen\textendash{}\allowbreak{}Kružík;
Ghorbel\textendash{}\allowbreak{}Launay & Abstract-level evidence and metadata; used for general methodological or
physical context \\
\hline
Ng; Sun; Mitsoulis\textendash{}\allowbreak{}Hatzikiriakos; Petit & Metadata/discovery excerpts; reading leads and broad subject scope, not
quantitative validation evidence \\
\hline
Yoon\textendash{}\allowbreak{}Lim & Partial publisher-text access through §5.4.2; complete limitations
unavailable \\
\bottomrule
\end{longtable}
\end{small}

The source-version qualifications apply to the technical descriptions
throughout this report. In particular, final publication metadata do not
prove that a preprint\textquotesingle s wording is identical to the
published version. Detailed locators and search queries are in
\texttt{evidence/\allowbreak{}visual\_\penalty 250 inverse.md},
\texttt{evidence/\allowbreak{}dough\_\penalty 250 rheology.md}, and
\texttt{evidence/\allowbreak{}methods\_\penalty 250 evaluation.md}. The
CSV catalogue provides primary links and access levels for each
reference.

\subsection{Bottom line}\label{bottom-line}

\textbf{The scientific opportunity is real, but it is an experimental
and inverse-problem question\textemdash{}\allowbreak{}not a novelty
claim about using Taichi or parameter sweeps.} A framework that
establishes when depth and tool kinematics can distinguish dough
response, quantifies what remains ambiguous, and predicts new
manipulations would be a credible contribution. The strongest report of
that work will connect numerical verification, measurement physics, food
preparation, and independent prediction in one transparent argument.

\section*{References}\label{bibliography}
\addcontentsline{toc}{section}{References}

\phantomsection\label{refs}
\begin{CSLReferences}{0}{0}
\bibitem[\citeproctext]{ref-review2020}
\CSLLeftMargin{{[}1{]} }%
\CSLRightInline{V. E. Arriola-Rios, P. Guler, F. Ficuciello, D. Kragic,
B. Siciliano, and J. L. Wyatt, {`Modeling of Deformable Objects for
Robotic Manipulation: A Tutorial and Review'}, \emph{Frontiers in
Robotics and AI}, vol. 7, Art. no. 82, 2020, doi:
\href{https://doi.org/10.3389/frobt.2020.00082}{10.3389/frobt.2020.00082}.}

\bibitem[\citeproctext]{ref-raue2011}
\CSLLeftMargin{{[}2{]} }%
\CSLRightInline{A. Raue, C. Kreutz, T. Maiwald, U. Klingmüller, and J.
Timmer, {`Addressing parameter identifiability by model-based
experimentation'}, \emph{IET Systems Biology}, vol. 5, no. 2, pp.
120\textendash{}\allowbreak{}130, 2011, doi:
\href{https://doi.org/10.1049/iet-syb.2010.0061}{10.1049/iet-syb.2010.0061}.}

\bibitem[\citeproctext]{ref-discrepancy2014}
\CSLLeftMargin{{[}3{]} }%
\CSLRightInline{J. Brynjarsdóttir and A. O'Hagan, {`Learning about
physical parameters: the importance of model discrepancy'},
\emph{Inverse Problems}, vol. 30, no. 11, Art. no. 114007, 2014, doi:
\href{https://doi.org/10.1088/0266-5611/30/11/114007}{10.1088/0266-5611/30/11/114007}.}

\bibitem[\citeproctext]{ref-kennedy2001}
\CSLLeftMargin{{[}4{]} }%
\CSLRightInline{M. C. Kennedy and A. O'Hagan, {`Bayesian calibration of
computer models'}, \emph{Journal of the Royal Statistical Society:
Series B (Statistical Methodology)}, vol. 63, no. 3, pp.
425\textendash{}\allowbreak{}464, 2001, doi:
\href{https://doi.org/10.1111/1467-9868.00294}{10.1111/1467-9868.00294}.}

\bibitem[\citeproctext]{ref-yazar2023}
\CSLLeftMargin{{[}5{]} }%
\CSLRightInline{G. Yazar, {`Wheat Flour Quality Assessment by
Fundamental Non-Linear Rheological Methods: A Critical Review'},
\emph{Foods}, vol. 12, no. 18, Art. no. 3353, 2023, doi:
\href{https://doi.org/10.3390/foods12183353}{10.3390/foods12183353}.}

\bibitem[\citeproctext]{ref-ng2006}
\CSLLeftMargin{{[}6{]} }%
\CSLRightInline{T. S. K. Ng, G. H. McKinley, and M. Padmanabhan,
{`Linear to Non-linear Rheology of Wheat Flour Dough'}, \emph{Applied
Rheology}, vol. 16, no. 5, pp. 265\textendash{}\allowbreak{}274, 2006,
doi:
\href{https://doi.org/10.1515/arh-2006-0019}{10.1515/arh-2006-0019}.}

\bibitem[\citeproctext]{ref-sun2020}
\CSLLeftMargin{{[}7{]} }%
\CSLRightInline{X. Sun, F. Koksel, M. T. Nickerson, and M. G. Scanlon,
{`Modeling the viscoelastic behavior of wheat flour dough prepared from
a wide range of formulations'}, \emph{Food Hydrocolloids}, vol. 98, Art.
no. 105129, 2020, doi:
\href{https://doi.org/10.1016/j.foodhyd.2019.05.030}{10.1016/j.foodhyd.2019.05.030}.}

\bibitem[\citeproctext]{ref-anderssen2013}
\CSLLeftMargin{{[}8{]} }%
\CSLRightInline{R. S. Anderssen and M. Kružík, {`Modelling of
wheat-flour dough mixing as an open-loop hysteretic process'},
\emph{Discrete and Continuous Dynamical Systems
\textendash{}\allowbreak{} B}, vol. 18, no. 2, pp.
283\textendash{}\allowbreak{}293, 2013, doi:
\href{https://doi.org/10.3934/dcdsb.2013.18.283}{10.3934/dcdsb.2013.18.283}.}

\bibitem[\citeproctext]{ref-rolling2009}
\CSLLeftMargin{{[}9{]} }%
\CSLRightInline{E. Mitsoulis and S. G. Hatzikiriakos, {`Rolling of bread
dough: Experiments and simulations'}, \emph{Food and Bioproducts
Processing}, vol. 87, no. 2, pp. 124\textendash{}\allowbreak{}138, 2009,
doi:
\href{https://doi.org/10.1016/j.fbp.2008.07.001}{10.1016/j.fbp.2008.07.001}.}

\bibitem[\citeproctext]{ref-cocuzza2018}
\CSLLeftMargin{{[}10{]} }%
\CSLRightInline{S. Cocuzza and X.-T. Yan, {`First engineering framework
for the out-of-plane robotic shaping of thin rheological objects'},
\emph{Robotics and Computer-Integrated Manufacturing}, vol. 53, pp.
108\textendash{}\allowbreak{}121, 2018, doi:
\href{https://doi.org/10.1016/j.rcim.2018.02.005}{10.1016/j.rcim.2018.02.005}.}

\bibitem[\citeproctext]{ref-adhesion2014}
\CSLLeftMargin{{[}11{]} }%
\CSLRightInline{D. Ghorbel and B. Launay, {`An investigation into the
nature of wheat flour dough adhesive behaviour'}, \emph{Food Research
International}, vol. 64, pp. 305\textendash{}\allowbreak{}313, 2014,
doi:
\href{https://doi.org/10.1016/j.foodres.2014.06.045}{10.1016/j.foodres.2014.06.045}.}

\bibitem[\citeproctext]{ref-fabbri2015}
\CSLLeftMargin{{[}12{]} }%
\CSLRightInline{A. Fabbri and C. Cevoli, {`Estimation of semolina dough
rheological parameters by inversion of a finite elements model'},
\emph{Journal of Agricultural Engineering}, vol. 46, no. 3, pp.
95\textendash{}\allowbreak{}99, 2015, doi:
\href{https://doi.org/10.4081/jae.2015.469}{10.4081/jae.2015.469}.}

\bibitem[\citeproctext]{ref-miranda2025}
\CSLLeftMargin{{[}13{]} }%
\CSLRightInline{A. Monleón-Getino \emph{et al.}, {`Advancing
Viscoelastic Material Characterization Through Computer Vision and
Robotics: MIRANDA and RELAPP'}, \emph{Materials}, vol. 18, no. 21, Art.
no. 4827, 2025, doi:
\href{https://doi.org/10.3390/ma18214827}{10.3390/ma18214827}.}

\bibitem[\citeproctext]{ref-pizza2017}
\CSLLeftMargin{{[}14{]} }%
\CSLRightInline{A. Petit, V. Lippiello, G. A. Fontanelli, and B.
Siciliano, {`Tracking elastic deformable objects with an RGB-D sensor
for a pizza chef robot'}, \emph{Robotics and Autonomous Systems}, vol.
88, pp. 187\textendash{}\allowbreak{}201, 2017, doi:
\href{https://doi.org/10.1016/j.robot.2016.08.023}{10.1016/j.robot.2016.08.023}.}

\bibitem[\citeproctext]{ref-dpsi2025}
\CSLLeftMargin{{[}15{]} }%
\CSLRightInline{X. Yang, Z. Ji, and Y.-K. Lai, {`Differentiable
physics-based system identification for robotic manipulation of
elastoplastic materials'}, \emph{The International Journal of Robotics
Research}, vol. 44, no. 13, pp. 2126\textendash{}\allowbreak{}2155,
2025, doi:
\href{https://doi.org/10.1177/02783649251334661}{10.1177/02783649251334661}.}

\bibitem[\citeproctext]{ref-empm2026}
\CSLLeftMargin{{[}16{]} }%
\CSLRightInline{Y. Chen, Y. Hu, L. Sun, T. Kusnur, L. Herlant, and C.
Jiang, {`EMPM: Embodied MPM for Modeling and Simulation of Deformable
Objects'}, \emph{IEEE Robotics and Automation Letters}, vol. 11, no. 4,
pp. 4179\textendash{}\allowbreak{}4186, 2026, doi:
\href{https://doi.org/10.1109/LRA.2026.3664610}{10.1109/LRA.2026.3664610}.}

\bibitem[\citeproctext]{ref-diffcal2023}
\CSLLeftMargin{{[}17{]} }%
\CSLRightInline{K. Arnavaz, M. K. Nielsen, P. G. Kry, M. Macklin, and K.
Erleben, {`Differentiable Depth for Real2Sim Calibration of Soft Body
Simulations'}, \emph{Computer Graphics Forum}, vol. 42, no. 1, pp.
277\textendash{}\allowbreak{}289, 2023, doi:
\href{https://doi.org/10.1111/cgf.14720}{10.1111/cgf.14720}.}

\bibitem[\citeproctext]{ref-diffcloud2022}
\CSLLeftMargin{{[}18{]} }%
\CSLRightInline{P. Sundaresan, R. Antonova, and J. Bohg, {`DiffCloud:
Real-to-Sim from Point Clouds with Differentiable Simulation and
Rendering of Deformable Objects'}, in \emph{IEEE/RSJ International
Conference on Intelligent Robots and Systems (IROS)}, 2022, pp.
10828\textendash{}\allowbreak{}10835. doi:
\href{https://doi.org/10.1109/IROS47612.2022.9981101}{10.1109/IROS47612.2022.9981101}.}

\bibitem[\citeproctext]{ref-wang2015}
\CSLLeftMargin{{[}19{]} }%
\CSLRightInline{B. Wang, L. Wu, K. Yin, U. Ascher, L. Liu, and H. Huang,
{`Deformation Capture and Modeling of Soft Objects'}, \emph{ACM
Transactions on Graphics}, vol. 34, no. 4, Art. no. 94, 2015, doi:
\href{https://doi.org/10.1145/2766911}{10.1145/2766911}.}

\bibitem[\citeproctext]{ref-hahn2019}
\CSLLeftMargin{{[}20{]} }%
\CSLRightInline{D. Hahn, P. Banzet, J. M. Bern, and S. Coros,
{`Real2Sim: Visco-elastic parameter estimation from dynamic motion'},
\emph{ACM Transactions on Graphics}, vol. 38, no. 6, Art. no. 236, 2019,
doi:
\href{https://doi.org/10.1145/3355089.3356548}{10.1145/3355089.3356548}.}

\bibitem[\citeproctext]{ref-pacnerf2023}
\CSLLeftMargin{{[}21{]} }%
\CSLRightInline{X. Li \emph{et al.}, {`PAC-NeRF: Physics Augmented
Continuum Neural Radiance Fields for Geometry-Agnostic System
Identification'}, in \emph{International Conference on Learning
Representations (ICLR)}, 2023. Available:
\url{https://openreview.net/forum?id=tVkrbkz42vc}}

\bibitem[\citeproctext]{ref-gic2024}
\CSLLeftMargin{{[}22{]} }%
\CSLRightInline{J. Cai \emph{et al.}, {`GIC: Gaussian-Informed Continuum
for Physical Property Identification and Simulation'}, in \emph{Advances
in Neural Information Processing Systems}, 2024, pp.
75035\textendash{}\allowbreak{}75063. doi:
\href{https://doi.org/10.52202/079017-2388}{10.52202/079017-2388}.}

\bibitem[\citeproctext]{ref-masiv2025}
\CSLLeftMargin{{[}23{]} }%
\CSLRightInline{Y. Zhao \emph{et al.}, {`Toward Material-Agnostic System
Identification from Videos'}, in \emph{Proceedings of the IEEE/CVF
International Conference on Computer Vision (ICCV)}, 2025, pp.
5944\textendash{}\allowbreak{}5956. doi:
\href{https://doi.org/10.1109/ICCV51701.2025.00562}{10.1109/ICCV51701.2025.00562}.}

\bibitem[\citeproctext]{ref-phystwin2025}
\CSLLeftMargin{{[}24{]} }%
\CSLRightInline{H. Jiang, H.-Y. Hsu, K. Zhang, H.-N. Yu, S. Wang, and Y.
Li, {`PhysTwin: Physics-Informed Reconstruction and Simulation of
Deformable Objects from Videos'}, in \emph{Proceedings of the IEEE/CVF
International Conference on Computer Vision (ICCV)}, 2025, pp.
7219\textendash{}\allowbreak{}7230. Available:
\url{https://openaccess.thecvf.com/content/ICCV2025/papers/Jiang_PhysTwin_Physics-Informed_Reconstruction_and_Simulation_of_Deformable_Objects_from_Videos_ICCV_2025_paper.pdf}}

\bibitem[\citeproctext]{ref-matl2020}
\CSLLeftMargin{{[}25{]} }%
\CSLRightInline{C. Matl, Y. Narang, R. Bajcsy, F. Ramos, and D. Fox,
{`Inferring the Material Properties of Granular Media for Robotic
Tasks'}, in \emph{IEEE International Conference on Robotics and
Automation (ICRA)}, 2020. doi:
\href{https://doi.org/10.1109/ICRA40945.2020.9197063}{10.1109/ICRA40945.2020.9197063}.}

\bibitem[\citeproctext]{ref-yoon2025}
\CSLLeftMargin{{[}26{]} }%
\CSLRightInline{K. Yoon and S.-C. Lim, {`Real-to-sim high-resolution
cloth modeling: Physical parameter optimization using particle-based
simulation with robot manipulation data'}, \emph{Journal of
Computational Design and Engineering}, vol. 12, no. 8, pp.
29\textendash{}\allowbreak{}44, 2025, doi:
\href{https://doi.org/10.1093/jcde/qwaf065}{10.1093/jcde/qwaf065}.}

\bibitem[\citeproctext]{ref-mls2018}
\CSLLeftMargin{{[}27{]} }%
\CSLRightInline{Y. Hu \emph{et al.}, {`A Moving Least Squares Material
Point Method with Displacement Discontinuity and Two-Way Rigid Body
Coupling'}, \emph{ACM Transactions on Graphics}, vol. 37, no. 4, Art.
no. 150, 2018, doi:
\href{https://doi.org/10.1145/3197517.3201293}{10.1145/3197517.3201293}.}

\bibitem[\citeproctext]{ref-apic2015}
\CSLLeftMargin{{[}28{]} }%
\CSLRightInline{C. Jiang, C. Schroeder, A. Selle, J. Teran, and A.
Stomakhin, {`The Affine Particle-In-Cell Method'}, \emph{ACM
Transactions on Graphics}, vol. 34, no. 4, Art. no. 51, 2015, doi:
\href{https://doi.org/10.1145/2766996}{10.1145/2766996}.}

\bibitem[\citeproctext]{ref-snow2013}
\CSLLeftMargin{{[}29{]} }%
\CSLRightInline{A. Stomakhin, C. Schroeder, L. Chai, J. Teran, and A.
Selle, {`A Material Point Method for Snow Simulation'}, \emph{ACM
Transactions on Graphics}, vol. 32, no. 4, Art. no. 102, 2013, doi:
\href{https://doi.org/10.1145/2461912.2461948}{10.1145/2461912.2461948}.}

\bibitem[\citeproctext]{ref-ils2020}
\CSLLeftMargin{{[}30{]} }%
\CSLRightInline{C. Liu and W. Sun, {`ILS-MPM: An implicit
level-set-based material point method for frictional particulate contact
mechanics of deformable particles'}, \emph{Computer Methods in Applied
Mechanics and Engineering}, vol. 369, Art. no. 113168, 2020, doi:
\href{https://doi.org/10.1016/j.cma.2020.113168}{10.1016/j.cma.2020.113168}.}

\bibitem[\citeproctext]{ref-chainqueen2019}
\CSLLeftMargin{{[}31{]} }%
\CSLRightInline{Y. Hu \emph{et al.}, {`ChainQueen: A Real-Time
Differentiable Physical Simulator for Soft Robotics'}, in \emph{IEEE
International Conference on Robotics and Automation (ICRA)}, 2019, pp.
6265\textendash{}\allowbreak{}6271. doi:
\href{https://doi.org/10.1109/ICRA.2019.8794333}{10.1109/ICRA.2019.8794333}.}

\bibitem[\citeproctext]{ref-difftaichi2020}
\CSLLeftMargin{{[}32{]} }%
\CSLRightInline{Y. Hu \emph{et al.}, {`DiffTaichi: Differentiable
Programming for Physical Simulation'}, in \emph{International Conference
on Learning Representations (ICLR)}, 2020. Available:
\url{https://openreview.net/forum?id=B1eB5xSFvr}}

\bibitem[\citeproctext]{ref-robocraft2022}
\CSLLeftMargin{{[}33{]} }%
\CSLRightInline{H. Shi, H. Xu, Z. Huang, Y. Li, and J. Wu, {`RoboCraft:
Learning to See, Simulate, and Shape Elasto-Plastic Objects with Graph
Networks'}, in \emph{Robotics: Science and Systems XVIII}, 2022. doi:
\href{https://doi.org/10.15607/RSS.2022.XVIII.008}{10.15607/RSS.2022.XVIII.008}.}

\bibitem[\citeproctext]{ref-robocook2023}
\CSLLeftMargin{{[}34{]} }%
\CSLRightInline{H. Shi, H. Xu, S. Clarke, Y. Li, and J. Wu, {`RoboCook:
Long-Horizon Elasto-Plastic Object Manipulation with Diverse Tools'}, in
\emph{Proceedings of the 7th Conference on Robot Learning}, 2023, pp.
642\textendash{}\allowbreak{}660. Available:
\url{https://proceedings.mlr.press/v229/shi23a.html}}

\bibitem[\citeproctext]{ref-plasticinelab2021}
\CSLLeftMargin{{[}35{]} }%
\CSLRightInline{Z. Huang \emph{et al.}, {`PlasticineLab: A Soft-Body
Manipulation Benchmark with Differentiable Physics'}, in
\emph{International Conference on Learning Representations (ICLR)},
2021. Available: \url{https://openreview.net/forum?id=xCcdBRQEDW}}

\bibitem[\citeproctext]{ref-doughnet2024}
\CSLLeftMargin{{[}36{]} }%
\CSLRightInline{D. Bauer, Z. Xu, and S. Song, {`DoughNet: A Visual
Predictive Model for Topological Manipulation of Deformable Objects'},
in \emph{Computer Vision \textendash{}\allowbreak{} ECCV 2024}, 2024,
pp. 92\textendash{}\allowbreak{}108. doi:
\href{https://doi.org/10.1007/978-3-031-72940-9_6}{10.1007/978-3-031-72940-9\_6}.}

\bibitem[\citeproctext]{ref-raue2009}
\CSLLeftMargin{{[}37{]} }%
\CSLRightInline{A. Raue \emph{et al.}, {`Structural and practical
identifiability analysis of partially observed dynamical models by
exploiting the profile likelihood'}, \emph{Bioinformatics}, vol. 25, no.
15, pp. 1923\textendash{}\allowbreak{}1929, 2009, doi:
\href{https://doi.org/10.1093/bioinformatics/btp358}{10.1093/bioinformatics/btp358}.}

\bibitem[\citeproctext]{ref-ral-scope}
\CSLLeftMargin{{[}38{]} }%
\CSLRightInline{IEEE Robotics and Automation Society, {`IEEE Robotics
and Automation Letters: Information for Authors'}. Accessed: Sep. 11,
2026. {[}Online{]}. Available:
\url{https://www.ieee-ras.org/publications/ra-l/ra-l-information-for-authors/}}

\bibitem[\citeproctext]{ref-rap-scope}
\CSLLeftMargin{{[}39{]} }%
\CSLRightInline{IEEE Robotics and Automation Society, {`IEEE Robotics
and Automation Practice'}. Accessed: Sep. 11, 2026. {[}Online{]}.
Available: \url{https://www.ieee-ras.org/publications/ra-p/}}

\bibitem[\citeproctext]{ref-joss-scope}
\CSLLeftMargin{{[}40{]} }%
\CSLRightInline{Journal of Open Source Software, {`Submitting a paper to
JOSS'}. Accessed: Sep. 11, 2026. {[}Online{]}. Available:
\url{https://joss.readthedocs.io/en/latest/submitting.html}}

\end{CSLReferences}

\end{document}
